\chapter{信息受限下的强化学习}
\label{chap:rl-information}

一辆小车驶近岔路口，摄像头清楚地拍到了它的位置，控制器仍可能无法决定怎样转向。如果画面没有显示小车向哪一侧滑动，控制器需要借助过去的观测辨认当前运动状态。如果位置和速度都已知，但日志从未记录某种转向动作，学习者需要判断现有经验能否说明该动作的后果。这两种困难需要不同的信息：记忆可以保留过去的线索，却不能补出从未尝试的行为；反复训练可以重新利用日志，却不能找回未被记录的线索。

第~\ref{chap:rl-learning}~章已经说明怎样从经验估计价值、优化策略和学习模型。本章考察这些方法在信息受限时还能使用到哪一步。状态不可见时，我们要从当前观测转向完整历史，或构造保留决策所需信息的历史表示。经验固定时，我们要先判断日志是否足以确定候选策略的价值，也就是价值是否可识别，然后才能讨论估计和改进。图~\ref{fig:rl-information-roadmap}~将这两条认识路径分开。离线部分可观测问题同时受到两类限制：既要检查历史中留下了哪些线索，也要检查收集者依据什么选择了日志中的动作。

\begin{figure}[htbp]
  \centering
  % Knowledge dependencies; final width: 112 mm.
\begin{tikzpicture}[
  x=\linewidth,y=1mm,
  font=\figurefont\fontsize{8}{10}\selectfont,
  box/.style={draw=BookRule,fill=BookPaper,line width=.8pt,
    text=BookInk,text width=.22\linewidth,minimum height=11mm,
    align=center,inner sep=1.5mm},
  flow/.style={-{Latex[length=1.8mm]},draw=BookInk,line width=.9pt}
]
  \node[box] (limits) at (.13,0) {两类信息限制};
  \node[box,fill=BookTeal!8] (history) at (.5,0) {历史与决策状态};
  \node[box,fill=BookBlue!8] (combined) at (.87,0) {叠加限制\\与可识别性};
  \node[box] (evaluation) at (.13,-26) {固定经验\\评价与选择};
  \node[box] (learning) at (.5,-26) {固定经验\\学习与改进};
  \node[box] (rounds) at (.87,-26) {有限交互\\修正};
  \draw[flow] (limits) -- (history);
  \draw[flow] (history) -- (combined);
  \draw[flow] (limits) -- (evaluation);
  \draw[flow] (evaluation) -- (learning);
  \draw[flow] (learning) -- (rounds);
  \draw[flow] (rounds) -- (combined);
  \draw[flow] (evaluation.north east) -- (.27,-13)
    -- (.82,-13) -- (combined.south west);
  \draw[flow] (learning.north) -- (.5,-16)
    -- (.82,-16) -- (combined.south);
\end{tikzpicture}

  \caption{本章的阅读依赖。箭头表示理解后项需要前项建立的概念，不表示算法的执行顺序。历史表示说明行动时能依据什么，固定数据评价说明已有经验能支持哪些判断；两者共同限制离线部分可观测问题中的策略结论。}
  \label{fig:rl-information-roadmap}
\end{figure}

% Outline: 4.1
\section{两类信息限制}
\label{sec:rl-information-limits}

% Outline: 4.1.1
\subsection{决策所依据的状态信息}
\label{sec:rl-information-state-information}

第~\ref{chap:rl-overview}~章的马尔可夫条件针对完整状态$s_t$：给定它和当前动作，就能描述下一步的条件分布，无须再查询更早的经历。摄像头取得的$o_t$只是状态的一部分信息，给定同一画面和动作，后果仍可能随未见的速度而变化。因此，执行策略不能直接沿用以$s_t$为输入的规则，而要考虑此前已取得的记录。一步记录$(o_t,a_t,r_t,o_{t+1})$按真实交互时刻$t$排列：先看到$o_t$，再执行$a_t$，随后取得奖励$r_t$和下一观测。这个时间下标不表示训练中的梯度更新次数。

\begin{definition}[决策历史]
  \label{def:rl-information-history}
  时刻$t$的完整记录历史按时间排列行动前已经取得的观测、动作和奖励，写为
  \[
    h_t=(o_0,a_0,r_0,\ldots,a_{t-1},r_{t-1},o_t).
  \]
  历史条件策略$\pi(a\mid h_t)$以行动前已获得的观测、动作和奖励为依据。当前观测策略的输入仅为$o_t$；长度为$L$的窗口策略仅保留最近$L$步记录，并须约定窗口在回合边界处怎样构造。
\end{definition}

历史中也保留奖励，因为奖励可能帮助辨认状态。例如，碰撞后返回的奖励可以说明小车触及了画面中看不见的障碍；删去它，就可能丢掉下一次决策需要的线索。这里保留的都是已经发生的反馈，未来才揭示的状态或奖励不能提前交给当前策略。

\begin{example}[相同位置与不同记忆]
  \label{ex:rl-information-cue}
  考虑一个独立教学任务。小车的真实状态包括位置$x_t$和速度$v_t$，摄像头只给出位置。同样看到$x_t=0$时，若$v_t=1$，所需制动方向与$v_t=-1$时相反。为判断运动方向，可以比较相邻位置。在无噪声、采样间隔为1且速度在相邻采样间不变的简化条件下，$v_t=x_t-x_{t-1}$；两次位置观测足以恢复速度，单帧没有这项信息。

  再考虑线索出现得更早的岔路任务。回合开始时，提示$c\in\{L,R\}$以相同概率显示；小车随后经过$L+1$步外观相同的走廊，途中只能前进，奖励均为0。出口处选中$c$指示的方向得1，选错得0。本例及后文均按不折扣的回合成功奖励比较，其期望就是成功概率，不采用默认折扣回报。完整历史策略可以记住提示，期望成功奖励达到1。最近$L$步窗口已经看不到提示，固定的前进动作也不能编码提示，因此这类策略的最好期望成功奖励为$1/2$。
\end{example}

位置任务只需多保留一帧，便能补出缺失的速度；岔路任务却说明，任意预先固定的窗口都可能因线索延迟更长而失效。表示是否足够，取决于它有没有保留区分动作后果的信息，而不是输入序列单纯有多长。

% Outline: 4.1.2
\subsection{学习能够访问的经验}
\label{sec:rl-information-data-access}

现在让日志完整记录小车的位置和速度。如果收集者在某个状态从不向左转，我们仍可能不知道在那里左转会怎样。缺口不在单条记录的字段，而在整批记录涉及的行为。分析可用经验时，必须分别检查记录是否完整、收集行为覆盖了哪里，以及学习者能否主动补采样。

\begin{definition}[经验访问协议]
  \label{def:rl-information-access}
  \term{固定经验}协议将可用材料限定为给定数据集$\mathcal D$及声明可用的元数据，训练和选模期间都不允许查询真实环境。\term{有限新增交互}协议允许在约定的轮次和预算内补充真实经验；每轮可执行的策略、重置权限和专家帮助权限均由协议规定。\term{持续在线交互}协议允许策略根据经验不断更新并继续收集数据，但重置、传感和测试权限仍须明确，不能据此假定可以从任意状态采样。
\end{definition}

固定日志即使回放一百万轮，真实访问过的状态和动作也不会增加，改变的只是算法利用这些记录的方式。学习模型生成一百万条模拟转移，同样没有新增真实查询。表~\ref{tab:rl-information-access}~沿用完整记录位置和速度的小车任务，对照交互权限改变后能够取得的证据。

\begin{table*}[htbp]
  \centering
  \small
  \begin{tabularx}{\linewidth}{p{20mm}XXXX}
    \toprule
    协议 & 记录字段 & 行为覆盖 & 主动采样权限 & 评价查询权限 \\
    \midrule
    固定经验 & 位置、速度、动作、奖励和回合边界 & 仅原日志的收集行为 & 无真实调用、无重置 & 只能离线估计；不能试跑候选 \\
    有限新增 & 原字段加策略版本、轮次和查询用途 & 原日志加预算内新增行为 & 每轮限定步数；重置单列计数 & 允许的测试也扣除预算 \\
    持续在线 & 同样记录状态与行为元数据 & 随收集策略持续变化 & 在可达轨迹上继续行动；任意状态重置仍需授权 & 独立测试须保留其数据职责与成本 \\
    \bottomrule
  \end{tabularx}
  \caption{同一完全可观测小车任务中的三种经验协议。完整字段不能补齐未尝试的动作；可以继续交互，也不意味着可以像调用生成模型那样从任意指定状态采样。}
  \label{tab:rl-information-access}
\end{table*}

% Outline: 4.1.3
\subsection{观测完整性与数据访问权的组合}
\label{sec:rl-information-two-axes}

小车能否取得完整状态，并不决定训练者能否继续驾驶它。只见位置的控制器可以被允许主动加减速，借反馈辨认运动方向；已经记录位置和速度的日志也可以禁止新增试验。表~\ref{tab:rl-information-two-axes}~将观测完整性与交互权限交叉排列，得到四种条件。

\begin{table}[htbp]
  \centering
  \small
  \begin{tabularx}{\linewidth}{p{18mm}XX}
    \toprule
    执行观测 & 可新增真实交互 & 只有固定经验 \\
    \midrule
    完整状态 & 位置、速度均可见，可继续控制小车 & 状态字段完整，但只能使用已有驾驶日志 \\
    部分观测 & 只见位置，可主动探查，构成在线POMDP & 只见位置的固定轨迹，构成离线POMDP \\
    \bottomrule
  \end{tabularx}
  \caption{观测完整性与数据访问权可以交叉组合。允许少量新增经验只放宽访问权限，不改变观测类型；表中的POMDP将在下节正式定义。}
  \label{tab:rl-information-two-axes}
\end{table}

下节先固定执行时不能直接看到状态的条件，考察怎样由已有观测形成决策依据。第~\ref{sec:rl-information-evaluation}、\ref{sec:rl-information-offline-learning}~节改回完全可观测任务，分别检查固定日志怎样支持评价和学习。第~\ref{sec:rl-information-combined}~节再将两类限制叠加，重新核对成立条件。分开讨论以后，才能辨清失败来自历史压缩遗漏线索，还是来自日志没有覆盖目标动作。

% Outline: 4.2
\section{决策信息不完整下的强化学习}
\label{sec:rl-information-partial-observation}

% Outline: 4.2.1
\subsection{从当前观测扩展为观测历史}
\label{sec:rl-information-history-policy}

回到例~\ref{ex:rl-information-cue}~的出口：同一画面下，早期提示为左时应向左，为右时应向右。\term{观测混叠}（observation aliasing）就是不同决策处境对应同一个可见输入，因而仅凭该输入无法区分它们。完整隐藏状态可以满足马尔可夫条件，观测过程却未必满足；我们需要同时描述状态怎样演化、观测怎样产生。图~\ref{fig:rl-information-aliasing}~将这两种关系分开，也显示相同当前观测可以来自不同历史。

\begin{definition}[有限部分可观测马尔可夫决策过程]
  \label{def:rl-information-pomdp}
  \term{部分可观测马尔可夫决策过程}（partially observable Markov decision process, POMDP）分别描述环境的隐藏状态和智能体取得的观测。本章采用有限状态集$\mathcal S$、动作集$\mathcal A$、观测集$\mathcal O$、初始分布$\rho_0$与有限时域$H$；为使用有限求和，暂将奖励取值也限定在有限集$\mathcal R$。

  初始状态$s_0\sim\rho_0$生成观测$o_0\sim O_0(\cdot\mid s_0)$。在时刻$t$，智能体依据$h_t$选择$a_t$，环境随后生成$s_{t+1}$及反馈$y_{t+1}=(r_t,o_{t+1})$。将下一状态与反馈的条件分布合写为已知联合核
  \[
    K(s',y\mid s,a)
    =\mathbb P(s_{t+1}=s',y_{t+1}=y\mid s_t=s,a_t=a).
  \]
  对反馈求和得到转移核$P(s'\mid s,a)=\sum_yK(s',y\mid s,a)$；按奖励值加权所有后果，得到期望奖励
  $R(s,a)=\sum_{s',r,o'}rK(s',(r,o')\mid s,a)$。奖励有界，采用回报
  $\sum_{t=0}^{H-1}\gamma^t r_t$，其中$0\leq\gamma<1$，终端价值取0。
\end{definition}

联合核保留了奖励与观测可能共享的隐藏信息，只有满足相应独立性时，才可进一步分解为转移和观测模型。与这一信息约定对应，执行策略写作$\pi_t(a\mid h_t)$；仅使用当前观测的$\pi_t(a\mid o_t)$限制了可选择的策略。我们保留时刻$t$，因为有限时域中剩余时间也会影响动作选择，即使对隐藏状态的推断相同也如此。

\begin{figure}[htbp]
  \centering
  % Two cue histories with identical final observations.
\begin{tikzpicture}[
  x=\linewidth,y=1mm,
  font=\figurefont\fontsize{8}{10}\selectfont,
  state/.style={draw=BookBlue,fill=BookBlue!7,line width=.8pt,
    minimum width=21mm,minimum height=8mm,align=center},
  obs/.style={text=BookInk,align=center},
  evolution/.style={-{Latex[length=1.8mm]},draw=BookInk,line width=.9pt},
  observation/.style={-{Latex[length=1.8mm]},draw=BookMuted,
    line width=.8pt,dashed}
]
  \node[state] (l0) at (.13,0) {$c=L$};
  \node[state] (lm) at (.47,0) {保留$c=L$};
  \node[state] (lt) at (.81,0) {到达出口};
  \node[obs] (ol0) at (.13,-13) {提示$L$};
  \node[obs] (olm) at (.47,-13) {相同走廊};
  \node[obs] (olt) at (.81,-13) {相同出口画面};
  \node[obs,text=BookTeal] at (.48,-23)
    {历史策略：记住提示$L$，在出口选左};
  \node[state] (r0) at (.13,-35) {$c=R$};
  \node[state] (rm) at (.47,-35) {保留$c=R$};
  \node[state] (rt) at (.81,-35) {到达出口};
  \node[obs] (or0) at (.13,-48) {提示$R$};
  \node[obs] (orm) at (.47,-48) {相同走廊};
  \node[obs] (ort) at (.81,-48) {相同出口画面};
  \node[obs,text=BookTeal] at (.48,-58)
    {历史策略：记住提示$R$，在出口选右};
  \draw[evolution] (l0) -- (lm);
  \draw[evolution] (lm) -- node[above] {$\cdots$} (lt);
  \draw[evolution] (r0) -- (rm);
  \draw[evolution] (rm) -- node[above] {$\cdots$} (rt);
  \foreach \a/\b in {l0/ol0,lm/olm,lt/olt,r0/or0,rm/orm,rt/ort}
    \draw[observation] (\a) -- (\b);
  \node[font=\figurefont\fontsize{7.5}{9}\selectfont,text=BookMuted]
    at (.48,-70) {实线：环境演化\qquad 虚线：观测生成};
\end{tikzpicture}

  \caption{早期提示任务中的两条可能轨迹。实线描述隐藏状态演化，虚线描述观测生成，途中固定为前进的动作未画出。出口画面相同，旁边单列的历史策略却因记住不同提示而选择相反方向；这些动作不属于观测。}
  \label{fig:rl-information-aliasing}
\end{figure}

仿真器有时能在训练日志中给出$s_t$，帮助训练评论家或表示。但部署时若只能取得$o_t$，执行策略就不能读取这一特权状态，评价时也须提供部署真正可用的历史。要利用状态模型而不直接观测状态，可以推断当前状态的概率分布；下面的信念方法正是通过这一分布连接推断与决策，经典POMDP论述也采用这一思路\scite{rl1-kaelbling1998-pomdp}

% Outline: 4.2.2
\subsection{将观测历史压缩为决策状态}
\label{sec:rl-information-sufficiency}

完整历史保留了已记录的线索，但输入和存储都随时间增长。压缩历史时，我们关心的是哪些信息会改变行动后果的预测，而不是能否复原每个过去细节。模型已知时，可以将历史转换为当前各个可能状态的概率。这样的概率分布称为信念状态，作用是保留当前状态的不确定性。

\begin{definition}[信念状态]
  \label{def:rl-information-belief}
  给定POMDP与完整记录历史，\term{信念状态}（belief state）定义为
  \[
    b_t(s)=\mathbb P(s_t=s\mid h_t),\qquad
    b_t\in\Delta(\mathcal S),
  \]
  其中$\Delta(\mathcal S)$是$\mathcal S$上的概率单纯形，即所有状态概率非负且总和为1的分布集合。$b_t$是给定历史后的后验分布，不表示另一个物理状态；只留下$\argmax_s b_t(s)$得到的点估计也不等于信念。
\end{definition}

推断从回合开始便已发生。对于具有正概率的初始观测$o_0$，按初始分布与观测似然加权并归一化，得到
$b_0(s)=\rho_0(s)O_0(o_0\mid s)/\sum_{\bar s}\rho_0(\bar s)O_0(o_0\mid\bar s)$；
仅当初始观测不提供状态信息时，$b_0$才可直接取为$\rho_0$。

如果把“障碍在左侧的概率为0.6”直接改成“障碍在左侧”，右侧仍有0.4概率这一信息就消失了。碰撞成本很高时，这部分概率可能改变应选的动作。要说明保留完整分布为何足以支持控制，需要证明：预测奖励、更新状态推断和比较动作，都可以只使用信念完成。

\begin{theorem}[精确信念对最优控制充分]
  \label{thm:rl-information-belief-sufficient}
  考虑定义~\ref{def:rl-information-pomdp}~中的已知有限POMDP，并假设动作选择仅依赖完整记录历史与独立的策略随机数。每个时刻$t$的最优历史价值都能写成$V_t^*(h_t)=U_t^*(b_t)$，且存在仅以$(t,b_t)$为依据的最优策略。
\end{theorem}

\begin{proof}
  先检查信念能否预测一个动作的后果。固定正概率历史$h_t$并选择$a$；由于动作只依据该历史和独立随机数，给定历史后的动作本身不再透露隐藏状态。按当前各状态的概率加权联合核，全概率公式给出
  \[
    \mathbb P(s_{t+1}=s',y_{t+1}=y\mid h_t,a)
    =\sum_s b_t(s)K(s',y\mid s,a).
  \]
  为预测可见反馈，将上式对隐藏的下一状态$s'$求和，得到
  $q(y\mid b_t,a)=\sum_{s,s'}b_t(s)K(s',y\mid s,a)$；将奖励取期望，得到
  $\bar R(b_t,a)=\sum_s b_t(s)R(s,a)$。反馈概率和奖励期望都由$(b_t,a)$确定。

  反馈出现后还需要更新推断。对$q(y\mid b_t,a)>0$的反馈，以其概率归一化状态与反馈的联合概率，贝叶斯公式给出
  \[
    b_{t+1}(s')
    =\frac{\sum_s b_t(s)K(s',y\mid s,a)}
           {q(y\mid b_t,a)}
    =:\tau(b_t,a,y)(s').
  \]
  新信念只需旧信念、动作和新反馈便可算出，这就是后验递推闭合；已经压缩的其余历史无需再读。零概率反馈不贡献期望，因此任意指定这些反馈上的$\tau$都不改变价值。

  最后用反向归纳检查最优动作比较。终点满足$V_H^*(h_H)=0$，可令$U_H^*(b)=0$。假设$t+1$时已有$V_{t+1}^*(h)=U_{t+1}^*(b(h))$；那么在$t$时选择$a$，当前奖励加上按可能反馈加权的最优继续价值为
  \[
    \bar R(b_t,a)
    +\gamma\sum_y q(y\mid b_t,a)
                      U_{t+1}^*(\tau(b_t,a,y)).
  \]
  上式仅依赖$b_t$与$a$。动作都来自有限动作集，按期望回报从中取最大值，便定义了$U_t^*(b_t)$并完成这一步归纳。在每个$(t,b)$处选取最大化动作，以固定规则处理平局，就构造出达到该值的信念策略。归纳到$t=0$，结论成立。
\end{proof}

这一结果说明最优控制不必保留信念以外的历史，并没有要求所有历史策略都在相同信念下选择同一动作。控制器可以故意按一个与未来无关的旧观测切换动作；这只是它使用了额外历史，不影响最优信念策略的存在。

上述构造需要已知模型；模型未知时，精确后验也不能直接计算。学习型潜在信念尝试用$\bm z_t$编码隐藏状态的不确定性。另一种选择是\term{预测状态表示}（predictive state representation, PSR），用若干未来试验发生的预测概率概括历史。“试验”包含指定的未来动作和随后要观察的反馈，因此只预测收集策略自然产生的下一帧还不够。若这些概率能够递推所有相关试验，就可不必为隐藏状态赋予语义名称；表示是否有限维、能否从数据学习，仍依赖结构条件\scite{rl4-littman2001-psr}

选择学习目标时，还要明确表示对什么问题充分。表~\ref{tab:rl-information-sufficiency}~分别考察未来反馈、固定策略价值和最优控制。例如，在两个隐藏状态中，“始终停止”的价值均为0，但向左、向右的奖励恰好相反。常数表示能准确预测停止策略的价值，却无法帮助选择转向动作。固定策略价值充分，不能据此推出改进所需的充分性。

\begin{table*}[htbp]
  \centering
  \small
  \begin{tabularx}{\linewidth}{p{25mm}XXX}
    \toprule
    充分性要求 & 必须保留的量 & 动作或策略范围 & 仍可丢失什么 \\
    \midrule
    未来反馈预测充分 & 各允许动作试验下的未来反馈联合分布与可递推状态 & 声明的全部试验，通常跨多个动作与时域 & 对这些预测无影响的历史细节与隐藏状态名称 \\
    给定策略价值充分 & 给定$\pi$和奖励下的条件期望回报 & 只要求该$\pi$的继续行为；还需说明其执行输入 & 其他动作的价值、完整反馈分布 \\
    最优控制充分 & 允许策略类中达到最优值所需的动作比较与更新 & 所允许的控制行为；依赖任务与时域 & 不影响最优动作的预测细节 \\
    \bottomrule
  \end{tabularx}
  \caption{三种充分性对应不同问题。固定策略的回报预测正确，仍可能缺少换策略或换奖励所需的信息；已知POMDP中的精确信念则具有最优控制充分性的依据。}
  \label{tab:rl-information-sufficiency}
\end{table*}

% Outline: 4.2.3
\subsection{新观测到来后的决策状态更新}
\label{sec:rl-information-filtering}

信念充分性证明中，每次取得反馈都要生成新信念。实际计算时，给定已执行的$a_t$和收到的$y_{t+1}=(r_t,o_{t+1})$，先遍历可能的前态，用旧信念加权联合核，得到
$\widetilde b_{t+1}(s')=\sum_s b_t(s)K(s',y_{t+1}\mid s,a_t)$。这是下一状态与该反馈同时出现的概率，尚未归一化。再对所有下一状态求和，算出模型赋予该反馈的总概率；以它为分母，得到给定反馈后的状态分布：
\begin{equation}
  \label{eq:rl-information-bayes}
  b_{t+1}(s')
  =\frac{\sum_s b_t(s)K(s',y_{t+1}\mid s,a_t)}
         {\sum_{\bar s}\sum_s b_t(s)K(\bar s,y_{t+1}\mid s,a_t)}.
\end{equation}
这个条件分布只在分母为正时有定义。真实反馈若在模型中概率为0，滤波器便无法按上式归一化，需要检查模型或传感假设是否与实际过程失配。

若希望先预测状态、再单独用观测校正，就需要对联合核增加条件。假设
$K(s',(r,o')\mid s,a)=P(s'\mid s,a)O(o'\mid s',a)p(r\mid o',a)$，
这表示给定$o'$与$a$后，奖励不再透露前后状态的信息。此时奖励项在分子分母中相同，可以约去；先对前态加权转移核，再按新观测的似然重新加权和归一化，得到
\begin{align}
  \label{eq:rl-information-predict-correct}
  b_{t+1}^{-}(s')&=\sum_sP(s'\mid s,a_t)b_t(s),\\
  b_{t+1}(s')&=
  \frac{O(o_{t+1}\mid s',a_t)b_{t+1}^{-}(s')}
       {\sum_{\bar s}O(o_{t+1}\mid\bar s,a_t)b_{t+1}^{-}(\bar s)}.
  \nonumber
\end{align}
碰撞奖励若由隐藏状态决定且提供了观测以外的线索，这一分解条件便不成立，更新时仍须保留奖励。

有限状态可以直接更新每个状态的概率；连续状态则需要适合其分布结构的表示。在线性高斯系统中，后验始终是高斯分布，只需存储均值和协方差。设
$\bm x_{t+1}=\bm A\bm x_t+\bm B\bm u_t+\bm w_t$，
$\bm o_{t+1}=\bm C\bm x_{t+1}+\bm v_{t+1}$；
其中初始状态为高斯，过程噪声与观测噪声是相互独立的零均值高斯量，协方差分别为$\bm W,\bm N$，并与初始状态独立；另行记录的奖励暂不提供额外状态信息。当前后验为
$\mathcal N(\bm m_t,\bm\Sigma_t)$时，Kalman滤波先传播均值和协方差，再用实际观测与预测观测的差修正它们。增益按预测不确定性和观测噪声决定修正幅度，计算为
\begin{align}
  \label{eq:rl-information-kalman}
  \bm m^-_{t+1}&=\bm A\bm m_t+\bm B\bm u_t,&
  \bm\Sigma^-_{t+1}&=\bm A\bm\Sigma_t\bm A^\top+\bm W,\nonumber\\
  \bm L_{t+1}&=\bm\Sigma^-_{t+1}\bm C^\top
     (\bm C\bm\Sigma^-_{t+1}\bm C^\top+\bm N)^{-1},\nonumber\\
  \bm m_{t+1}&=\bm m^-_{t+1}
       +\bm L_{t+1}(\bm o_{t+1}-\bm C\bm m^-_{t+1}),\nonumber\\
  \bm\Sigma_{t+1}&=(\bm I-\bm L_{t+1}\bm C)\bm\Sigma^-_{t+1}.
\end{align}
这里要求创新协方差可逆；创新就是实际观测相对预测观测的差。增益$\bm L_{t+1}$按两种不确定性的大小分配预测与新观测的权重。递推之所以精确，是因为线性高斯条件使后验保持高斯，而非均值和协方差能够代表任意分布\scite{rl4-kalman1960-filter}

\begin{example}[一次标量滤波]
  \label{ex:rl-information-kalman}
  用标量教学系统$x_{t+1}=x_t+u_t+w_t$、$o_{t+1}=x_{t+1}+v_{t+1}$核算一次更新。当前$m_t=0$、$\Sigma_t=1$，令$u_t=1$、$W=N=1$，随后观测到2。状态预测为$m^-_{t+1}=1$、$\Sigma^-_{t+1}=2$，由此算得增益$2/3$。观测比预测大1，校正得到$m_{t+1}=1+(2/3)(2-1)=5/3$、$\Sigma_{t+1}=2/3$。方差尚未归零，所以$5/3$是后验均值，不是已知的真实位置。
\end{example}

显式模型若不可信，就不能直接依靠上述精确滤波。可以改为从数据学习递推记忆
$\bm z_{t+1}=f_{\bm\theta}(\bm z_t,a_t,r_t,o_{t+1})$，
也可以保留历史，再由注意力网络重新计算$\bm z_t$。图~\ref{fig:rl-information-memory}~对照了两种计算方式。输入序列必须保持同一轨迹的时间关系，记忆按协议在回合开始重置；任意训练片段的开头未必是回合开始，不能在那里直接清零。

\begin{figure*}[htbp]
  \centering
  % Recursive and attention interfaces; final width: 169 mm.
\begin{tikzpicture}[
  x=\linewidth,y=1mm,
  font=\figurefont\fontsize{8}{10}\selectfont,
  box/.style={draw=BookRule,fill=BookPaper,line width=.8pt,
    minimum height=9mm,align=center,inner sep=2mm},
  flow/.style={-{Latex[length=1.8mm]},draw=BookInk,line width=.9pt},
  guide/.style={draw=BookMuted,line width=.8pt,dashed}
]
  \node[anchor=west,text=BookInk] at (0,12) {(a) 递推记忆};
  \node[anchor=west,text=BookInk] at (.55,12) {(b) 历史注意力};
  \node[box] (old) at (.055,0) {$\bm z_t$};
  \node[box,fill=BookTeal!8] (update) at (.23,0) {$f_{\bm\theta}$};
  \node[box] (new) at (.405,0) {$\bm z_{t+1}$};
  \node[box] (feedback) at (.23,-17) {$a_t,r_t,o_{t+1}$};
  \node[box] (policy) at (.405,-17) {策略};
  \draw[flow] (old) -- (update);
  \draw[flow] (update) -- (new);
  \draw[flow] (feedback) -- (update);
  \draw[flow] (new) -- (policy);
  \node[box] (history) at (.59,0) {$h_{t+1}$};
  \node[box,fill=BookTeal!8] (attention) at (.77,0) {注意力};
  \node[box] (encoding) at (.945,0) {$\bm z_{t+1}$};
  \node[box] (append) at (.655,-17) {追加动作与反馈};
  \node[box] (apolicy) at (.945,-17) {策略};
  \draw[flow] (history) -- (attention);
  \draw[flow] (attention) -- (encoding);
  \draw[flow] (append) -- (history);
  \draw[flow] (encoding) -- (apolicy);
  \draw[guide] (.01,-33) rectangle (.3,-44);
  \draw[guide] (.3,-33) rectangle (.99,-44);
  \node at (.155,-38.5) {预热前缀：重建记忆};
  \node at (.645,-38.5) {训练后缀：计算损失};
  \draw[guide] (.01,-29) -- (.01,-48);
  \draw[guide] (.3,-29) -- (.3,-48);
  \node[anchor=north west] at (.01,-49) {真实回合开始：重置};
  \node[anchor=north west] at (.32,-49) {片段损失起点：不盲目清零};
  \draw[guide] (.23,-33) -- (feedback.south);
  \draw[guide] (.655,-33) -- (append.south);
\end{tikzpicture}

  \caption{递推记忆与历史注意力的输入和更新。实线描述表示与策略的计算，虚线标出训练片段及边界。预热前缀用于恢复损失起点处的记忆，不计算其训练损失。图示片段从真实回合开始；回合中部片段需用已保存状态或更早前缀初始化，片段边界本身不构成重置依据。}
  \label{fig:rl-information-memory}
\end{figure*}

DRQN将循环状态加入价值网络，使前帧信息可以传到后帧。相对于堆叠有限帧，它改变了历史的保存方式，价值目标仍采用时序差分学习\scite{rl4-hausknecht2015-drqn}
序列进入重放后，还要处理参数变化导致的表示漂移和旧隐藏状态，R2D2分析了这些问题。截断反向传播只限制梯度跨越多少步，并未说明更早的历史无用。burn-in即预热，先运行片段前缀恢复隐藏状态，再在后缀计算损失。也可以保存收集时的隐藏状态，减少从零启动的偏差；但网络参数更新后，保存的状态可能已不适合当前网络\scite{rl4-kapturowski2019-r2d2}

注意力模型也可以跨片段保存信息，GTrXL通过门控等结构调整适配强化学习优化。外部记忆则将历史写入独立存储，再按需读取。增加这些保存能力会带来缓存、读写和训练成本。更长的上下文或更低的预测损失，只描述输入范围或拟合结果，尚不能证明表示满足定理~\ref{thm:rl-information-belief-sufficient}~的条件。

% Outline: 4.2.4
\subsection{基于决策状态的策略求解与动作选择}
\label{sec:rl-information-belief-control}

信念更新给出当前各状态的概率，我们还需要据此决定动作。比较时要平均可能状态的后果，并允许后续动作随新反馈改变，不能先挑一个最可能状态当成真相。\term{信念MDP}（belief MDP）把可计算的信念本身作为完全可见的决策状态。其即时奖励为
$\bar R(b,a)=\sum_s b(s)R(s,a)$，反馈概率为
$q(y\mid b,a)=\sum_{s,s'}b(s)K(s',y\mid s,a)$，状态转移通过$\tau(b,a,y)$实现。
要计算剩余时域的最优回报，可以对有限动作集$\mathcal A$中的每个动作，加上即时奖励和按反馈概率加权的继续价值，再选最大的期望回报。第~\ref{chap:rl-planning}~章的Bellman关系由此写成
\begin{equation}
  \label{eq:rl-information-belief-bellman}
  U_t^*(b)=\max_a\left\{
    \bar R(b,a)+\gamma\sum_yq(y\mid b,a)
                    U_{t+1}^*(\tau(b,a,y))\right\},
  \qquad U_H^*=0.
\end{equation}
反馈不可由控制器挑选，必须逐个更新信念、计算继续价值后按发生概率取期望。一般不能用
$U_{t+1}^*(\mathbb E[b_{t+1}])$替代$\mathbb E[U_{t+1}^*(b_{t+1})]$，因为后续动作依据实际反馈而非平均反馈选择；先平均信念会丢掉这种适应机会。

\begin{example}[Tiger中的倾听、更新与开门]
  \label{ex:rl-information-tiger}
  Tiger任务有两扇门，一侧藏着危险，另一侧可以获得奖励。隐藏状态$s\in\{L,R\}$表示危险的方向。打开安全门得10，打开危险门得$-100$，随后回合终止；倾听得$-1$，状态不变，以0.95的概率报告正确方向。令初始$b_0(L)=b_0(R)=1/2$，最多行动两次，$\gamma=0.9$，时域结束后价值为0。这些数值与两步终止规则仅用于教学，不代表所有Tiger版本。

  先比较不获取信息的动作。直接打开任意一扇门，按初始危险概率加权，期望奖励都是
  $(1/2)10+(1/2)(-100)=-45$。若先倾听并得到“左侧有声”，用正确报告和错误报告的概率更新左侧危险概率，得到
  \[
    b_1(L)=\frac{0.95(1/2)}
                        {0.95(1/2)+0.05(1/2)}=0.95.
  \]
  按更新后的信念，开右门的期望奖励为$0.95(10)+0.05(-100)=4.5$，
  开左门为$0.95(-100)+0.05(10)=-94.5$；
  再次倾听只能得$-1$，之后已无动作可用，因此在这三个动作中应开右门。
  若听见右侧，按对称计算应开左门。两种反馈各以$1/2$发生，将各自的最优开门价值加权、折扣，再扣除倾听成本，初始倾听价值为
  \[
    Q_0(b_0,\text{倾听})
       =-1+0.9\bigl[(1/2)4.5+(1/2)4.5\bigr]=3.05.
  \]
  在此两步设定下，先倾听再按反馈开门的期望回报为3.05，高于直接开门的$-45$。
\end{example}

图~\ref{fig:rl-information-tiger}~显示了按反馈分开的动作比较。若固定同一个开门时刻，并把信号设为免费，最优开门收益的期望增量为$4.5-(-45)=49.5$；实际两步方案还包含倾听成本与后续折扣，不能直接使用这个增量。模型和听觉准确率虽然已知，当前危险位置仍未知，所以倾听可能有用。剩余时间不足、先验已很确定或探测成本过高时，获得信息的收益可能不足以支付代价，直接开门反而更合适。

\begin{figure}[htbp]
  \centering
  % Teaching parameters: accuracy .95, cost -1, discount .9, H=2.
\begin{tikzpicture}[
  x=\linewidth,y=1mm,
  font=\figurefont\fontsize{8}{10}\selectfont,
  box/.style={draw=BookRule,fill=BookPaper,line width=.8pt,
    align=center,minimum height=10mm,inner sep=2mm},
  flow/.style={-{Latex[length=1.8mm]},draw=BookInk,line width=.9pt}
]
  \node[box] (prior) at (.45,0) {$b_0(L)=0.5$};
  \node[box,fill=BookAmber!8] (open) at (.14,-22)
    {直接开门\\期望$-45$};
  \node[box,fill=BookTeal!8] (listen) at (.64,-22)
    {倾听：成本$-1$\\两步价值$3.05$};
  \node[box] (left) at (.43,-46)
    {听见左侧\\$b_1(L)=0.95$};
  \node[box] (right) at (.85,-46)
    {听见右侧\\$b_1(L)=0.05$};
  \node[box] (openright) at (.43,-68)
    {开右门\\期望$4.5$};
  \node[box] (openleft) at (.85,-68)
    {开左门\\期望$4.5$};
  \draw[flow] (prior) -- (open);
  \draw[flow] (prior) -- (listen);
  \draw[flow] (listen) -- node[above left] {$1/2$} (left);
  \draw[flow] (listen) -- node[above right] {$1/2$} (right);
  \draw[flow] (left) -- (openright);
  \draw[flow] (right) -- (openleft);
\end{tikzpicture}

  \caption{教学Tiger的两步决策树。$b(L)$是左门后有危险的概率。倾听反馈各以$1/2$发生，随后按新信念选择安全概率较高的门。末端标注条件期望开门奖励，根节点的3.05则计入倾听成本和折扣。}
  \label{fig:rl-information-tiger}
\end{figure}

若只保留先验概率最大的状态，平局时任取一侧为危险，再照完全可观测解打开另一侧，预测奖励会是10，真实期望却是$-45$。这不只是平局处理的问题。将先验改为$b_0(L)=0.6$，点估计仍直接开右门，但真实期望为$0.6(10)+0.4(-100)=-34$；按0.95准确率计算两个倾听分支，先听再开的价值仍是3.05。被删除的另一侧概率影响了信息是否值得获取，也就影响了当前动作。

信念足以支持最优控制，不意味着信念空间规划便宜。有限状态控制器用有限记忆节点规定动作分布和反馈后的跳转，限制了能实现的历史策略。点基值迭代保留信念的概率含义，仅在选定信念上近似价值备份。两者分别近似策略记忆和价值函数，不能混作同一种压缩；点基方法中采样了哪些信念，还会影响规划质量\scite{rl4-pineau2003-pbvi}

在特定控制问题中，估计与控制确实可以分开。线性、高斯、模型已知且没有额外约束的二次成本控制满足相应条件时，LQG分离原则允许把Kalman估计的状态均值交给线性最优反馈控制器。噪声和观测机制还须与动作无关；无限时域稳定解另需可稳定性、可检测性等条件。动作若能改变观测质量，或额外奖励透露了新的状态信息，就不能直接沿用该结论。Tiger中的倾听恰好改变信息获取，因此一般POMDP没有先独立估计、再按已知状态控制的普遍保证。

模型未知时，可以把学习到的$\bm z_t$接入第~\ref{chap:rl-learning}~章的方法。循环价值学习用它构造自举目标，循环行动者—评论家用它计算动作分布和优势，潜在模型则在该表示上学习转移与奖励，用于规划或想象轨迹。这些方法学习表示、价值或动力学，但还须检查表示是否能判断未实际执行的其他动作会怎样，即是否保留了反事实动作所需的信息。

% Outline: 4.2.5
\subsection{近似决策状态的误差与计算边界}
\label{sec:rl-information-memory-errors}

选择历史表示时，除网络参数外，还要计算保存历史、更新表示和随后规划的成本。表~\ref{tab:rl-information-memory-cost}~按$n=|\mathcal S|$、历史长度$t$、窗口长度$L$、单步输入维度$d$、循环隐藏维度$m$和控制器节点数$k$列出直接实现的复杂度。这些是表中指定计算方式的成本，不是对全部模型的复杂度断言。

\begin{table*}[htbp]
  \centering
  \small
  \begin{tabularx}{\linewidth}{p{21mm}XXXX}
    \toprule
    表示 & 历史范围与存储 & 一次更新代价 & 决策或规划空间 & 充分性依据 \\
    \midrule
    完整历史 & 全部；$O(td)$ & 追加记录$O(d)$；重新读取至少随$t$增长 & 分支随时域增长的历史树 & 不主动丢弃已记录信息 \\
    精确信念 & 历史递推压缩；$O(n)$ & 稠密有限核$O(n^2)$ & $n-1$维单纯形；有限层可达点也可指数增长 & 已知正确模型下的充分性定理 \\
    有限窗口 & 最近$L$步；$O(Ld)$ & 缓冲追加$O(d)$，编码另计 & 窗口输入；未必马尔可夫 & 需证明任务只依赖该窗口 \\
    有限状态控制器 & 当前节点$O(1)$；节点表另计 & 表查找可为$O(1)$ & $k$个记忆节点上的策略 & 策略类限制；一般非充分 \\
    学习型记忆 & 循环$O(m)$；注意力缓存随$L$增长 & 稠密循环$O(m^2+md)$；标准全窗口注意力$O(L^2m+Lm^2)$ & 学习表示上的价值或策略 & 任务与模型假设，加独立验证 \\
    \bottomrule
  \end{tabularx}
  \caption{信息表示的成本与依据。注意力复杂度按宽度$m$的一层完整窗口计算，增量缓存可改变单步成本；有限状态控制器若按节点与反馈跳转，跳转表需$O(k^2|\mathcal Y|)$存储，若另依赖实际动作，还需乘$|\mathcal A|$。存在充分统计量不等于易于规划。}
  \label{tab:rl-information-memory-cost}
\end{table*}

例如，GTrXL中的门控与残差调整主要帮助训练保持稳定，注意力缓存和计算仍随实际保留的序列长度变化。能否较顺利地训练，与学到的表示能否充当精确决策状态，需要分别核查\scite{rl4-parisotto2020-gtrxl}

要理解压缩怎样影响动作，可以沿表示、初始化、模型和求解逐一追踪。两个历史若被编码为同一$\bm z$，却在同一动作下具有不同的反馈分布或价值，单凭$\bm z$就无法区分它们；增加训练轮数不能恢复这种表示遗漏。表示本来能保存提示，却在每个重放片段开始被清零，则是初始化错误。联合核估计不准和正确目标下优化尚未收敛，又是不同误差，各自需要对应的修正。

若要把这些误差转换为动作损失，就需要对动作价值的误差作出明确假设。固定历史$h$，在同一允许动作集上比较真实最优继续价值和近似值；假设各动作都有
$|\widehat Q(h,a)-Q^*(h,a)|\leq\delta$。分别按近似值和真实值选取最大化动作
$\hat a\in\argmax_a\widehat Q(h,a)$、$a^*\in\argmax_aQ^*(h,a)$；插入这两个动作的近似价值，真实损失满足
\[
  Q^*(h,a^*)-Q^*(h,\hat a)
  \leq
  [Q^*(h,a^*)-\widehat Q(h,a^*)]
  +[\widehat Q(h,\hat a)-Q^*(h,\hat a)]
  \leq2\delta.
\]
中间省去的差为非正，因为$\hat a$最大化近似值，即$\widehat Q(h,a^*)\leq\widehat Q(h,\hat a)$。这一界要求逐动作误差受控，平均重建损失没有直接提供该条件；没有表示误差到价值误差的论证，就不能把训练损失当作$\delta$。

例~\ref{ex:rl-information-cue}~还可以直接检验重建与控制的区别。沿用不折扣成功奖励，两个记忆节点分别保存左右提示即可达到期望值1，丢弃提示则至多达到$1/2$。即便窗口内每帧都完美重建，$1/2$的控制差距仍在。若奖励后来改为依赖另一条早期线索，仅保存$c$又可能不足，当前任务充分性不能扩展为对所有奖励目标充分。第~\ref{chap:rl-learning}~章讨论更一般的样本复杂度与函数类条件；本例先表明，表示丢掉的信息不能靠更多样本或更充分的优化补回。

% Outline: 4.3
\section{固定经验下的策略评价与选择}
\label{sec:rl-information-evaluation}

% Outline: 4.3.1
\subsection{固定经验与行为策略}
\label{sec:rl-information-logging}

下面暂时恢复完全可观测条件：每条记录都有完整状态，但学习者不能再到环境中试验。要用这些记录评价另一策略，先要弄清动作由谁、按什么概率选出。同一份状态动作表可能来自固定策略、多个策略的混合或持续更新的控制器；收集方式不同，概率校正和误差估计所用的条件也不同。

\begin{definition}[行为策略与目标策略]
  \label{def:rl-information-behavior}
  沿用前文的行为策略$\mu$与目标策略$\pi$：前者生成日志动作，后者是待评价或部署的决策规律。将有限时域中收集的$n$个回合写为
  \[
    \mathcal D=\{\tau_i\}_{i=1}^n,\qquad
    \tau_i=(s_{i,0},a_{i,0},r_{i,0},\ldots,s_{i,H}).
  \]
  本节采用有限时域价值
  $J_H(\pi)=\mathbb E_\pi[\sum_{t=0}^{H-1}\gamma^t r_t]$。
  若回合提前终止，则以吸收状态和零奖励补到$H$。$J_H$只累计有限时域，不等于无限时域的$J$；后文切换目标时将另作说明。
\end{definition}

单一马尔可夫行为策略可逐步记录$\mu(a_t\mid s_t)$。多个子策略混合时，必须区别抽取时机：每步重新按已知概率选择子策略，可以用该步混合权重加权动作概率；回合开始只抽一次，则历史会改变对所选子策略的后验，通常不能一直使用固定混合权重。后一种收集方式需要保存策略身份及相应概率并作适当的条件采样分析，或记录真正的$\mu_t(a_t\mid h_t)$；这里的历史包含行动前记录的完整状态、动作和奖励。身份字段说明了收集机制，却不能补齐某个条件层中缺失的目标动作。若策略随收集过程自适应更新，还须保存版本和选择机制。

这些动作概率最好来自当时实际使用的随机化机制。机制记录正确时，概率是已知量；事后拟合的$\widehat\mu$则含有模型误差。确定性收集者没有尝试的动作，即使密度网络给出非零概率，也仍没有行为支持。表~\ref{tab:rl-information-fields}~说明哪些日志字段支撑评价中的具体计算。

\begin{table}[htbp]
  \centering
  \small
  \begin{tabularx}{\linewidth}{p{29mm}X}
    \toprule
    字段 & 支撑的计算与省略后的风险 \\
    \midrule
    状态、动作、奖励、下一状态 & 构造回报、Bellman目标及目标策略动作概率 \\
    真终止标记 & 真正终止后不再自举；终止奖励仍记在$r_t$ \\
    时间截断标记 & 区分任务结束与记录停止；持续任务截断处通常仍需继续价值 \\
    回合标识与时刻 & 恢复轨迹概率比、记忆重置和回合级划分 \\
    行为概率及策略版本 & 确定动作选择机制，区分已知概率与估计概率 \\
    初始条件、任务与目标 & 说明评价的$\rho_0$、奖励定义及任务分布是否匹配 \\
    \bottomrule
  \end{tabularx}
  \caption{完整状态日志仍需要采样元数据。有限时域的真正$H$步终止与持续任务在$H$处被时间截断，具有不同的自举目标。}
  \label{tab:rl-information-fields}
\end{table}

恢复回合边界后，还要确定统计重复的单位。独立收集的$n$个回合可以作为独立评价单位，回合内相邻转移却不能直接当作独立重复。只有一条相关长轨迹时，分块估计或相关数据推断需要混合性等条件，即随时间间隔增加，相关性受到相应控制。把相邻转移随机拆到训练与测试中，可能泄露轨迹信息，也会高估独立样本数。

D4RL包含不同质量、收集者及混合任务的数据，因此报告时要保留具体版本、奖励定义和归一化分数的参照；仅称“专家混合数据”还不能确定实验输入\scite{rl4-fu2020-d4rl}
除了数据组成，RL Unplugged还区分任务域与选模协议。调参是否使用在线策略测试尤其需要声明：如果依据真实环境分数选择超参数，整个选模过程就已超出纯固定经验评价\scite{rl4-gulcehre2020-unplugged}

Minari按回合提供数据与元数据接口，但具体数据集的终止、截断、环境版本和收集者说明仍须核对，存储格式并不保证行为概率可得\scite{rl4-farama-minari}
对于OGBench这样的离线目标条件强化学习基准，还要报告无奖励数据的组成、目标怎样选择、奖励怎样重标记以及测试目标协议。目标条件改变后，成功率不能直接横比。OGBench已有ICLR 2025正式版本；这里用它说明协议需包含哪些内容，不据其分数判断方法普遍优劣\scite{rl4-park2025-ogbench}

% Outline: 4.3.2
\subsection{策略价值的可识别性与覆盖条件}
\label{sec:rl-information-coverage}

固定日志可以交给多个候选策略评价，但每个候选会使用不同的状态动作。比较分数之前，先假设能够无限重复原收集方式，检查是否仍会有不同环境给出同样日志、却赋予候选不同价值。这是\term{可识别性}（identifiability）问题：可观察的数据规律能否唯一确定目标量。若不能，增加同类样本也无法只靠估计精度解决。

\begin{definition}[策略价值可识别]
  \label{def:rl-information-identifiability}
  策略价值可识别性以给定的数据生成协议、可观察字段、环境结构假设、初始分布和目标策略为条件。若所有产生相同可观测数据总体分布的允许环境都给出相同目标策略价值，则该价值可识别。这是总体分布的性质，不说明有限样本估计已经准确。
\end{definition}

\begin{theorem}[固定经验不能识别未覆盖动作的价值]
  \label{thm:rl-information-unseen-action}
  对单状态、单步、两个动作的任务，若日志行为只执行其中一个动作，且没有将两动作奖励联系起来的结构假设，则另一动作的价值及其相对行为策略的优劣一般不可识别。
\end{theorem}

\begin{proof}
  用两个仅在未见动作上不同的环境检验定义。唯一非终止状态为$s$，执行$a_0$或$a_1$后均终止。在$M_+$与$M_-$中，$a_0$的奖励都为0，$a_1$的奖励分别为$+1$和$-1$；两环境的行为均为$\mu(a_0\mid s)=1$。

  无论日志多大，其中只有$(s,a_0,0,s_{\mathrm{term}})$，所以可观测分布完全相同。但目标$\pi(a_1\mid s)=1$的价值分别是$+1$和$-1$，行为基线价值均为0。只使用日志的程序收到相同输入，其输出分布也相同，无法在两个环境中都正确确定目标价值或排序。因此，按定义，这两项结论都不可识别。
\end{proof}

给$a_1$产生数值预测还不够。神经网络可以外推，结构化动力学也可以将$a_0$的信息迁移到$a_1$，但需要可核验的正确性假设说明两动作怎样联系，才可能排除上述一个环境。网络能表达这种联系，不表示日志已经证明它成立。

为检查哪些候选避开了这一缺口，可以比较候选诱导的轨迹分布和行为分布。候选可能使用而行为绝不产生的轨迹，无法通过日志重加权补出；即使没有这种绝对缺口，也要考察候选是否把大量概率放在日志极少涉及的区域。

\begin{definition}[支持、集中程度与部分覆盖]
  \label{def:rl-information-support}
  目标轨迹分布$p_\pi$相对行为轨迹分布$p_\mu$\term{绝对连续}，记$p_\pi\ll p_\mu$，表示行为概率为0的任何轨迹事件也具有零目标概率。在共同环境核的有限模型中，目标可达历史上的每个目标动作都须有行为支持。

  若将评价改为无限时域折扣占用而非有限整条轨迹，则用前文归一化的$d^\pi(s,a)$表示目标访问概率，以$\nu(s,a)$表示数据转移的归一化采样分布。集中程度系数衡量目标访问相对数据采样可能集中到多大程度，用本质上确界记录占用比的上界而非平均大小，即取遍$\nu$支持中的状态动作、忽略$\nu$零测集。其定义为
  \[
    C_\pi=\operatorname*{ess\,sup}_{(s,a)\sim\nu}
               \frac{d^\pi(s,a)}{\nu(s,a)},
  \]
  并要求$d^\pi\ll\nu$。部分覆盖表示上述条件仅在某些策略或状态动作区域成立。可达性只说明某动作序列在环境中可能到达一处，不说明收集者赋予它正概率。
\end{definition}

非零支持排除了绝对缺口，却没有控制有限样本中的困难。目标经常使用的动作若只有$10^{-6}$行为概率，小日志很可能一次都未收集到；即使出现，概率校正也可能赋予很大权重。估计器不同，所需控制量可能是逐时占用比、轨迹密度比或其矩，不能把$C_\pi$作为所有分析的同一个条件。

\begin{table*}[htbp]
  \centering
  \small
  \begin{tabularx}{\linewidth}{p{26mm}XX}
    \toprule
    希望得到的结论 & 需要支持的行为范围 & 额外结构与不能省略的检查 \\
    \midrule
    给定策略评价 & 该策略的相关轨迹或占用 & 正确奖励、初始分布与行为机制；模型外推须有依据 \\
    候选排序 & 候选共同模型下的性能差 & 不一定要分别点识别两个价值；差的符号必须有证据 \\
    相对基线改进 & 候选与基线的比较行为 & 选择后仍有效的误差控制；缺数据处可限制回到基线 \\
    策略类中最优 & 足以排除所有有竞争力的替代行为 & 通常需更广覆盖或更强结构；高回报轨迹本身不够 \\
    \bottomrule
  \end{tabularx}
  \caption{结论越强，通常需要检查越多行为。结构假设可改变所需覆盖，但必须说明它如何排除与日志一致的不同环境。}
  \label{tab:rl-information-claims}
\end{table*}

表~\ref{tab:rl-information-claims}~将覆盖要求对应到结论强度。日志可能完整显示一条奖励为10的路径，却没尝试岔路处的另一动作；后者通向奖励11或$-100$，都可与日志相容。已观察路径可行有数据依据，路径最优或偏离它的策略更好却没有同样依据。因此，少量高回报轨迹不能替代备选行为的覆盖。

% Outline: 4.3.3
\subsection{离线策略价值的估计}
\label{sec:rl-information-estimators}

价值可识别后，仍要由有限日志计算估计。\term{离策略评价}（off-policy evaluation, OPE）用其他行为策略收集的数据估计目标策略价值；这里还限制为不能新增交互。先固定待评价的$\pi$，假设回合独立、初始分布及环境核共同、行为动作概率已知，且$p_\pi\ll p_\mu$。沿用定理\ref{thm:rl-learning-importance-sampling}的矩条件，目标下每步奖励绝对可积。为确定日志应怎样加权，先按生成顺序展开有限时域马尔可夫策略的轨迹密度：
\[
  p_\pi(\tau)=\rho_0(s_0)
       \prod_{t=0}^{H-1}
       \pi_t(a_t\mid s_t)p(s_{t+1},r_t\mid s_t,a_t).
\]
将$\pi$换为$\mu$得到行为密度，再计算目标密度除以行为密度。共同初始项和环境项约去，只需将沿途的目标动作概率除以行为概率并相乘：
\begin{equation}
  \label{eq:rl-information-weights}
  w_{i,0:t}=\prod_{k=0}^{t}
       \frac{\pi_k(a_{i,k}\mid s_{i,k})}
            {\mu_k(a_{i,k}\mid s_{i,k})}.
\end{equation}
这一比率把行为日志中的期望转换为目标策略下的期望，复用了第~\ref{chap:rl-learning}~章的换测度关系。初始分布若改变却未校正，或行为实际使用了未记录信息，便不能按上式消去相应因子。

令$G_i=\sum_{t=0}^{H-1}\gamma^t r_{i,t}$。可以将每条完整回报乘以轨迹比，再对独立回合取平均；也可以对每个奖励只校正此前的动作，随后累计和平均。普通轨迹重要性采样与逐决策形式分别为
\begin{align}
  \label{eq:rl-information-is}
  \widehat J_{\mathrm{IS}}
    &=\frac1n\sum_iw_{i,0:H-1}G_i,\nonumber\\
  \widehat J_{\mathrm{PDIS}}
    &=\frac1n\sum_i\sum_{t=0}^{H-1}
                      \gamma^t w_{i,0:t}r_{i,t}.
\end{align}
奖励$r_t$在$a_t$执行后产生，校正其分布须包括$a_t$，无需包括之后的动作。截到该奖励的前缀密度比正是$w_{i,0:t}$，由此得到逐决策形式；不必另外要求目标概率为0的行为历史上后续权重条件均值为1。概率正确、支持和上述矩条件成立时，两式都无偏，但相乘的权重可能随时域快速增大。

\begin{example}[三条两步轨迹的概率校正]
  \label{ex:rl-information-ope-data}
  为核算概率校正，给定一个教学MDP：从$s_0$出发，任意动作都到$s_1$，再行动一步终止。两个时刻均可选$A,B$，奖励为$r_0(A)=1$、$r_0(B)=0$、$r_1(A)=2$、$r_1(B)=0$，奖励与转移都确定，$\gamma=0.9$。行为每步均匀选择，目标每步以$3/4$概率选择$A$。对每步目标奖励取期望再折扣，真实价值为
  $J_2(\pi)=3/4+0.9(2\times3/4)=2.1$。

  假设独立收集到$AA,BA,BB$三条轨迹。按上述动作概率计算前缀比和回报，得到
  \[
  \begin{array}{c|rrrr}
    \text{轨迹}&G_i&w_{i,0:0}&w_{i,0:1}
                    &\sum_t\gamma^t w_{i,0:t}r_{i,t}\\
    \hline
    AA&2.8&1.5&2.25&5.55\\
    BA&1.8&0.5&0.75&1.35\\
    BB&0&0.5&0.25&0
  \end{array}
  \]
  先乘完整轨迹比再平均，普通IS为$(2.25\times2.8+0.75\times1.8)/3=2.55$；
  逐奖励使用前缀比，逐决策IS为$(5.55+1.35)/3=2.3$。本批样本的两个估计都偏离2.1；无偏性说的是反复收集整批日志后的估计期望，不保证每批数值准确。
\end{example}

轨迹权重过大时，可以改动汇总方式以减轻少数样本的影响，但也会改变统计性质。自归一化估计用
$\widehat J_{\mathrm{SN}}=\sum_iw_iG_i/\sum_iw_i$，
用总权重而非$n$归一化，要求分母非零；截断估计则预先选择正阈值$c$，限制每条权重不超过$c$：
$\widehat J_{\mathrm{clip}}=n^{-1}\sum_i\min(w_i,c)G_i$。
例~\ref{ex:rl-information-ope-data}~取$w_i=w_{i,0:1}$时，
自归一化结果为$7.65/3.25\approx2.3538$；取$c=1$的未归一化截断结果为
$(2.8+0.75\times1.8)/3\approx1.3833$。前者使用随机分母，后者修改精确换测度权重，一般都会引入有限样本偏差。以拟合行为概率替代已知概率，也改变无偏性条件。有效样本量诊断
$(\sum_iw_i)^2/\sum_iw_i^2$检查总权重是否主要由少数轨迹承担，却不检验奖励或行为模型是否正确，因而不能证明价值估计正确。

若日志能支持目标继续动作，也可以不乘整条轨迹的概率比，改为拟合目标策略的Bellman关系。\term{拟合Q评价}（fitted Q evaluation, FQE）固定$\pi$，用记录奖励与下一步目标价值构造回归目标，因此只评价、不取动作最大值。有限时域中令$\widehat Q_H=0$，从末步向前更新。时刻$t$在允许函数类$\mathcal F_t$中选取$f$，使数据上的平方残差最小；下一动作按固定$\pi_{t+1}$取期望：
\begin{equation}
  \label{eq:rl-information-fqe}
  \widehat Q_t
  \in\argmin_{f\in\mathcal F_t}
    \sum_{i}
    \left[
      f(s_{i,t},a_{i,t})
      -r_{i,t}
      -\gamma\mathbb E_{a'\sim\pi_{t+1}(\cdot\mid s_{i,t+1})}
                       \widehat Q_{t+1}(s_{i,t+1},a')
    \right]^2.
\end{equation}
其中$\mathcal F_t$限定时刻$t$可拟合的价值函数。末步$t=H-1$的继续价值直接为0，无需定义终点策略。若改为无限折扣任务，则用时间不变的函数类反复回归，以上一轮$\widehat Q^{(k)}$构造目标，得到$\widehat Q^{(k+1)}$。两种形式都在真终止处屏蔽后续项。回归后，对初始状态及目标初始动作取期望，有限时域得到$\widehat J_H(\pi)=\mathbb E_{s_0\sim\rho_0,a_0\sim\pi_0}\widehat Q_0(s_0,a_0)$；初始分布未知时，需要相应初态样本。每轮拟合要扫描或采样数据、计算目标动作期望并求解回归，非线性函数类还会产生优化误差。小训练平方误差仅说明数据上的拟合，有限样本价值界另需函数类、Bellman逼近和覆盖条件\scite{rl4-le2019-fqe}

\begin{example}[在同一日志上完成FQE与DR]
  \label{ex:rl-information-fqe-dr}
  在例~\ref{ex:rl-information-ope-data}~的同三条轨迹上，每个时刻使用独立价值表。先拟合末步奖励，得到
  $\widehat Q_1(s_1,A)=2$、$\widehat Q_1(s_1,B)=0$，故
  $\widehat V_1(s_1)=0.75(2)+0.25(0)=1.5$。再将这个目标继续价值传回初始步，两个回归目标为
  \[
    \widehat Q_0(s_0,A)=1+0.9(1.5)=2.35,\qquad
    \widehat Q_0(s_0,B)=0+0.9(1.5)=1.35.
  \]
  按目标的初始动作概率加权，得到
  $\widehat J_{\mathrm{FQE}}=0.75(2.35)+0.25(1.35)=2.1$。
  本次结果精确，因为所需状态动作全已出现且后果确定；不能由此认为三条轨迹总是够用。若两时刻从零同时更新，第一轮只能拟合即时奖励，第二轮才会把1.5的继续价值传回初态。

  再将辅助预测与概率加权残差合用，观察预测不准时怎样校正。这种方法称为\term{双重稳健估计}（doubly robust estimation, DR），即权重或辅助价值有一侧正确时，在相应条件下仍可得到正确均值；具体条件将在一般公式后说明。为核算残差，事先给定不准确的辅助价值：
  $\widehat Q_1(A)=1$、$\widehat Q_1(B)=0$，
  $\widehat Q_0(A)=1.5$、$\widehat Q_0(B)=0.5$。
  对目标动作加权，得$\widehat V_1=0.75$、$\widehat V_0=1.25$。
  用实际奖励加预测继续价值减去当前预测，初始步残差$r_0+0.9\widehat V_1-\widehat Q_0$对两动作都为0.175，末步对$A$为1、对$B$为0。将这些残差按前缀比校正，加到初始预测上，三条轨迹在下面双重稳健公式中的贡献为
  \[
  \begin{aligned}
    D_{AA}&=1.25+1.5(0.175)+0.9(2.25)=3.5375,\\
    D_{BA}&=1.25+0.5(0.175)+0.9(0.75)=2.0125,\\
    D_{BB}&=1.25+0.5(0.175)=1.3375.
  \end{aligned}
  \]
  平均贡献为$2.2958$，尚有样本误差；若使用上面精确的FQE辅助值，各步残差均为0，每条轨迹都贡献2.1。
\end{example}

上例的DR使用有限时域逐决策形式，从预测初始价值出发，用实际反馈相对预测的Bellman残差修正。先按固定策略动作取期望，构造辅助
$\widehat V_t(s)=\mathbb E_{a\sim\pi_t}\widehat Q_t(s,a)$，
并令$\widehat V_H=0$；每个残差用截至该反馈的概率比转到目标分布，折扣累加后再平均回合贡献：
\begin{equation}
  \label{eq:rl-information-dr}
  \widehat J_{\mathrm{DR}}
  =\frac1n\sum_i\left[
    \widehat V_0(s_{i,0})
    +\sum_{t=0}^{H-1}\gamma^t w_{i,0:t}
       \bigl(r_{i,t}+\gamma\widehat V_{t+1}(s_{i,t+1})
                         -\widehat Q_t(s_{i,t},a_{i,t})\bigr)
    \right].
\end{equation}
“双重稳健”指权重或辅助价值有一侧正确时，在相应条件下仍可得到正确均值。为陈述精确无偏性，辅助函数须由独立训练数据拟合，或交叉拟合，使评价轨迹不参与自己的辅助训练。权重正确时，残差转到目标分布，$\widehat V$和$\widehat Q$项沿时间望远镜消去，留下$J_H(\pi)$，辅助值有误也不影响等式。辅助值全部满足$\widehat Q_t=Q_t^\pi$时，给定行动前历史与当前动作，残差条件均值为0，初始价值期望则正确；有限且在反馈前确定的错误行为权重也不改变均值。两侧都错便没有这种保证，任一侧正确也不保证方差小\scite{rl4-jiang2016-dr}

不使用行为概率的FQE仍需要覆盖。若上例在$s_1$只记录$B$，目标却选择$A$，$\widehat Q_1(s_1,A)$便无直接监督；其偏差$e$会以$\gamma\pi(A\mid s_1)e$进入初态目标，即使初态数据拟合完全正确也如此。模型评价遇到同样缺口时，传播的是奖励与转移预测误差。

轨迹比之外，还可以直接校正状态动作的分布。这里从有限时域切换到无限折扣、时间不变MDP，使用前文归一化的$d^\pi$。DICE类方法不必显式拟合$\mu$，而由占用流量平衡学习
$w(s,a)=d^\pi(s,a)/\nu(s,a)$；DualDICE通过对偶优化求解这类分布校正问题\scite{rl4-nachum2019-dualdice}
为确定比率，要求它对适当测试函数$f$恢复目标占用的流量关系：当前测试值减去折扣后的下一测试值，其加权期望等于初始输入质量。得到正确比率后再校正奖励期望：
\begin{align}
  \label{eq:rl-information-dice}
  &\mathbb E_{\nu,P,\pi}\!
       \left[w(s,a)\{f(s,a)-\gamma f(s',a')\}\right]
       =(1-\gamma)\mathbb E_{\rho_0,\pi}f(s_0,a_0),\nonumber\\
  &J(\pi)=\frac{1}{1-\gamma}\mathbb E_\nu[w(s,a)R(s,a)].
\end{align}
式中$a'\sim\pi(\cdot\mid s')$，条件正确的奖励样本可以替代奖励均值。$d^\pi$与$\nu$均归一化，而$J$是未归一化折扣回报，必须保留$(1-\gamma)^{-1}$。占用比校正仍要求目标受日志支持；省略行为拟合不会使支持外的价值可识别。

\begin{table*}[htbp]
  \centering
  \small
  \begin{tabularx}{\linewidth}{p{17mm}XXXX}
    \toprule
    估计路线 & 必需信息 & 拟合对象 & 关键覆盖要求 & 主要误差 \\
    \midrule
    IS & 轨迹、目标与行为概率 & 可不拟合模型 & 目标轨迹相对行为绝对连续 & 乘积权重方差；概率估计偏差 \\
    模型评价 & 转移、奖励及初始条件 & $\widehat P,\widehat R$ & 目标行为后果有支持或正确结构 & 模型误设、长滚动误差 \\
    FQE & 转移、固定$\pi$、初态 & $Q^\pi$ & 目标继续动作及状态可被评价 & 回归、投影、自举传播 \\
    本节DR & IS字段和辅助价值 & $\widehat Q,\widehat V$ & 仍需目标支持；辅助拟合与评价分离 & 双方错误的交互、权重方差 \\
    DualDICE & 转移分布、初态、$\pi$ & 占用比及对偶函数 & $d^\pi\ll\nu$与可控比值 & 比率函数逼近、优化与采样误差 \\
    \bottomrule
  \end{tabularx}
  \caption{固定策略的评价路线。表中列出的是各具体形式的条件，不将某一DR或DICE变种的结论推广到全部同名方法。}
  \label{tab:rl-information-estimators}
\end{table*}

表~\ref{tab:rl-information-estimators}~按同一固定策略任务汇总了所需信息和误差来源。评价器首先要与日志字段、支持条件相容，之后才比较方差、拟合成本和偏差控制；估计分数较高本身不是选择它的理由。

% Outline: 4.3.4
\subsection{候选策略的比较与选择}
\label{sec:rl-information-selection}

此前的估计针对固定$\pi$；若在同一日志上评价许多候选，再挑分数最高者，评价对象便由数据选择了。即便每个预固定候选都无偏，最大分数也容易选中正向噪声。候选和查询次数越多，直接报告被选中的最大值越可能乐观。

为控制选择影响，可以分清训练、选择与最终报告的数据职责。训练集产生候选和辅助模型，选择集用于调参或保留候选，最终评价集仅承担预先确定的报告。样本少时可采用严谨交叉拟合或统一误差控制；反复查看过的最终集已经参与选择，不能再称为未使用测试集。所有划分仍只重用原日志，不能创造缺失支持。

对于预先固定的两个候选，可以复用同一评价轨迹计算性能差。令$Z_i(\pi)$为IS或满足相应条件的DR轨迹贡献，先在每个回合算$\pi_1$贡献减去$\pi_2$贡献，再平均：
\begin{equation}
  \label{eq:rl-information-paired}
  \widehat\Delta=\frac1n\sum_i[Z_i(\pi_1)-Z_i(\pi_2)].
\end{equation}
这里估计的是$\pi_1$价值减去$\pi_2$价值，真实差为正表示$\pi_1$回报更高；有限样本的估计正值还要结合误差判断。固定候选与辅助模型、使用独立同分布评价回合且贡献的二阶矩有限时，为判断差的估计波动，将相减贡献的方差展开为
\[
  \operatorname{Var}(\widehat\Delta)
  =\frac1n\left[
    \operatorname{Var}(Z(\pi_1))+\operatorname{Var}(Z(\pi_2))
    -2\operatorname{Cov}(Z(\pi_1),Z(\pi_2))\right].
\]
共享轨迹可能使两个估计同向波动，正协方差会减小差的方差，负协方差则增大它。只看两个单独价值区间是否重叠，会遗漏比较所需的相关信息。

\begin{example}[价值区间重叠仍可比较性能差]
  \label{ex:rl-information-paired}
  构造两项价值估计5和4，各自标准误为1，两均值估计的协方差为0.98。相减得到性能差1，按方差公式计算标准误为
  $\sqrt{1+1-2(0.98)}=0.2$。另假设正态近似适用，边际95\%区间为$[3.04,6.96]$与$[2.04,5.96]$，高度重叠；配对差的95\%区间为$[0.608,1.392]$，却不包含0。这些统计量用于教学，不是实际运行实验所得。
\end{example}

图~\ref{fig:rl-information-comparison}~将单独价值与配对差的区间分开。此处价值在总体上可确定，区间来自有限样本误差；第~\ref{sec:rl-information-partial-identification}~节将改变这一条件，即使总体分布已知也可能存在多个相容价值。那里仍需比较共同模型下的差，但模型歧义与此处的采样波动是两类不确定性。

\begin{figure}[htbp]
  \centering
  % Teaching Gaussian intervals: marginal SE=1, covariance=.98.
\begin{tikzpicture}
  \begin{axis}[
    width=\linewidth,height=43mm,at={(0,0)},anchor=north west,
    xmin=0,xmax=8,ymin=.5,ymax=2.5,
    xtick={0,2,4,6,8},ytick={1,2},yticklabels={$\pi_2$,$\pi_1$},
    axis lines=left,xlabel={价值估计（回报）},
    tick label style={font=\figurefont\fontsize{8}{10}\selectfont},
    label style={font=\figurefont\fontsize{8}{10}\selectfont},
    axis line style={BookMuted,line width=.7pt},
    tick style={BookMuted},
    clip=false
  ]
    \addplot+[only marks,mark=*,mark size=2pt,BookTeal,
      error bars/.cd,x dir=both,x explicit,
      error bar style={line width=1pt}]
      coordinates {(5,2) +- (1.96,0)};
    \addplot+[only marks,mark=square*,mark size=2pt,BookBlue,
      error bars/.cd,x dir=both,x explicit,
      error bar style={line width=1pt}]
      coordinates {(4,1) +- (1.96,0)};
  \end{axis}
  \begin{axis}[
    width=\linewidth,height=36mm,at={(0,-49mm)},anchor=north west,
    xmin=-.2,xmax=1.6,ymin=.5,ymax=1.5,
    xtick={0,.4,.8,1.2,1.6},ytick={1},yticklabels={$\Delta$},
    axis lines=left,xlabel={性能差估计（回报）},
    tick label style={font=\figurefont\fontsize{8}{10}\selectfont},
    label style={font=\figurefont\fontsize{8}{10}\selectfont},
    axis line style={BookMuted,line width=.7pt},
    tick style={BookMuted},
    clip=false
  ]
    \addplot[BookMuted,dashed,line width=.8pt]
      coordinates {(0,.55) (0,1.45)};
    \addplot+[only marks,mark=diamond*,mark size=2.5pt,BookAmber,
      error bars/.cd,x dir=both,x explicit,
      error bar style={line width=1pt}]
      coordinates {(1,1) +- (.392,0)};
  \end{axis}
\end{tikzpicture}

  \caption{例~\ref{ex:rl-information-paired}~的教学区间。上组为两个价值的边际95\%正态近似区间，下组为共享轨迹计算的性能差区间；两组横轴各自标明量纲为回报。区间重叠与差是否包含0并非同一个判断。}
  \label{fig:rl-information-comparison}
\end{figure}

比较报告须列明候选数、产生候选所用数据、重复查询次数、置信水平和多重比较处理。候选集预先固定且有限时，可以分配各候选的失败概率，用并集界控制所有区间共同有效的事件；自适应增加候选不能直接沿用单候选区间。证据不足以排序时，就保留尚未排除的候选或报告无法排序。更多小数位不会增加证据。

% Outline: 4.4
\section{固定经验下的策略学习与改进}
\label{sec:rl-information-offline-learning}

% Outline: 4.4.1
\subsection{策略学习的目标与数据边界}
\label{sec:rl-information-learning-boundary}

评价固定策略只检查指定行为，学习则要由日志生成新的候选。可以先以记录行为为拟合目标。\term{行为克隆}（behavioral cloning, BC）就是让策略在日志状态上给记录动作较高概率；在参数化策略类中，通过降低数据动作的平均负对数似然拟合参数：
$\min_{\bm\phi}\mathbb E_{\mathcal D}[-\log\pi_{\bm\phi}(a\mid s)]$。
这一目标不使用奖励来重新分配概率。离线强化学习希望利用同一日志中的奖励选择行为，因此还要规定改进的比较基准。

拟合行为检查策略是否复现数据动作，提高估计回报检查学习目标是否上升。相对基线改进比较$J(\pi)$与$J(\pi_b)$，逼近最优则比较候选与允许策略类中最好的真实价值。后三者没有自动联系：优化器可能提高估计目标，却因查询数据外动作而利用模型错误，降低真实回报。

\begin{example}[一次贪心更新利用了外推错误]
  \label{ex:rl-information-extrapolation}
  在单步教学任务中，日志只记录奖励为1的$a_0$；未记录的$a_1$真实奖励为$-10$。网络预测$\widehat Q(s,a_0)=1$、$\widehat Q(s,a_1)=5$，数据动作平方误差为0。按两个动作中的最大预测值作一次贪心更新，会选中$a_1$：估计回报从1升为5，真实回报从1降为$-10$。反复拟合原日志中的$a_0$仍不会提供$a_1$的正确后果。
\end{example}

图~\ref{fig:rl-information-offline-pipeline}~把这类外推分成三个入口：策略生成新动作，价值备份查询新动作，规划器生成未观测状态动作。后续方法按限制这三个位置展开，也可组合使用。比较时保持数据与评价协议相同，说明优化对象、部署动作怎样产生以及额外假设；采用Transformer、扩散或流模型，并不单独说明策略是否改进。

\begin{figure*}[htbp]
  \centering
  % Fitting, extraction and evidence are different operations.
\begin{tikzpicture}[
  x=.97\linewidth,y=1mm,
  font=\figurefont\fontsize{8}{10}\selectfont,
  box/.style={draw=BookRule,fill=BookPaper,line width=.8pt,
    text width=.16\linewidth,minimum height=11mm,align=center,inner sep=1.5mm},
  flow/.style={-{Latex[length=1.8mm]},draw=BookInk,line width=.9pt},
  query/.style={-{Latex[length=1.8mm]},draw=BookAmber,line width=.8pt,dashed},
  lab/.style={fill=white,inner sep=1pt,font=\figurefont\fontsize{7.5}{9}\selectfont}
]
  \node[box] (data) at (.09,0) {固定数据\\$\mathcal D$};
  \node[box,fill=BookTeal!8] (q) at (.37,23) {价值模型\\$Q,V$};
  \node[box,fill=BookTeal!8] (behavior) at (.37,0) {行为模型\\$\widehat\mu$};
  \node[box,fill=BookTeal!8] (model) at (.37,-23) {环境模型\\$\widehat P,\widehat R$};
  \node[box] (extract) at (.65,0) {动作提取\\执行策略$\pi$};
  \node[box,fill=BookBlue!8] (evidence) at (.91,0) {评价证据\\保留日志或\\授权测试};
  \node[box] (objective) at (.09,23) {奖励目标\\与限制};
  \draw[flow] (objective) -- node[lab,above] {定义损失} (q);
  \draw[flow] (data.north east) -- node[lab,left] {拟合} (q.south west);
  \draw[flow] (data) -- node[lab,above] {拟合} (behavior);
  \draw[flow] (data.south east) -- node[lab,left] {拟合} (model.north west);
  \draw[flow] (q.east) -- (.65,23)
    -- node[lab,right] {价值引导} (extract.north);
  \draw[flow] (behavior) -- node[lab,above] {生成动作} (extract);
  \draw[flow] (model.east) -- (.65,-23)
    -- node[lab,right] {预测与规划} (extract.south);
  \draw[flow] (extract) -- node[lab,above] {评价} (evidence);
  \draw[query] (extract.north east) -- (.76,37)
    -- node[lab,above] {未覆盖动作：价值查询入口} (.37,37) -- (q.north);
  \draw[query] (extract.south east) -- (.79,-15) -- (.79,-39)
    -- node[lab,below] {未覆盖轨迹：模型外推入口} (.37,-39) -- (model.south);
\end{tikzpicture}

  \caption{固定数据上的学习与评价。实线分别标注拟合、生成动作及评价，虚线标出外推可能进入的查询。行为约束限制动作来源，悲观价值限制对未支持动作的乐观估计，模型约束限制规划范围；任何一条学习路径仍需独立的评价证据。}
  \label{fig:rl-information-offline-pipeline}
\end{figure*}

% Outline: 4.4.2
\subsection{行为数据约束下的策略学习}
\label{sec:rl-information-behavior-constraint}

在例~\ref{ex:rl-information-extrapolation}~中，问题由策略选择未见动作触发，可以先约束候选偏离记录行为的程度。以下保持完全可观测与固定$\mathcal D$，将价值目标改为无限折扣回报。TD3+BC沿用TD3的双评论家、目标网络和延迟行动者更新，在确定性行动者$f_{\bm\phi}(s)$的参数类中比较预测价值与动作平方偏差。对日志状态动作取平均，最大化价值减去偏差：
\begin{equation}
  \label{eq:rl-information-td3bc}
  \max_{\bm\phi}\ %
  \mathbb E_{(s,a)\sim\mathcal D}
  \left[\lambda Q_{\bm\theta}(s,f_{\bm\phi}(s))
            -\lVert f_{\bm\phi}(s)-a\rVert_2^2\right].
\end{equation}
动作是向量时，范数比较各动作坐标，此处沿用动作符号$a$。为调节价值项与距离项的相对尺度，在含$N$条记录的一批数据上常取
$\lambda=\alpha/(\frac1N\sum_i|Q_{\bm\theta}(s_i,a_i)|)$，
其中$\alpha$控制相对强度，行动者更新时固定$\lambda$，分母为0时实现须有非零保护。按链式法则对行动者参数求梯度，可以看到两项各自推动动作的方向：
\[
  \mathbb E_{\mathcal D}
  \left[(\nabla_{\bm\phi}f_{\bm\phi}(s))^\top
  \{\lambda\nabla_a Q_{\bm\theta}(s,a)|_{a=f_{\bm\phi}(s)}
          -2(f_{\bm\phi}(s)-a)\}\right].
\]
价值梯度推动动作朝预测价值较高处移动，距离梯度将它拉回记录动作。约束太弱可能再次利用外推，太强则接近模仿。两项的相对作用随奖励和动作尺度改变，因此比较$\alpha$必须同时给出归一化方式\scite{rl4-fujimoto2021-td3bc}

若日志行为有多个动作模式，平方拟合可能得到模式之间数据稀少的均值动作。可以让行为模型生成不同模式，再在候选附近改进。\term{批量约束Q学习}（batch-constrained Q-learning, BCQ）以这种方式缩小价值筛选的动作范围。连续动作版本先用$(s,a)$训练条件变分自编码器：编码器生成潜变量后验，解码器$g_{\bm\omega}(s,\bm z)$用后验样本重构动作；生成候选时改从规定先验采样，不读取待生成的数据动作。随后由$\xi_{\bm\phi}(s,a)$给每个坐标加入$[-\Phi,\Phi]$内的扰动，$\Phi\geq0$，并截到合法动作域。生成$m$个候选后，仅按$Q_1$在这些修正候选中取最大值，得到执行动作：
\begin{equation}
  \label{eq:rl-information-bcq}
  a_j=g_{\bm\omega}(s,\bm z_j),\qquad
  \widetilde a_j=\operatorname{clip}_{\mathcal A}
                  (a_j+\xi_{\bm\phi}(s,a_j)),\qquad
  a_{\mathrm{exec}}\in\argmax_{\widetilde a_j,\ 1\leq j\leq m}
                  Q_1(s,\widetilde a_j).
\end{equation}
行为生成器以动作重构误差加潜变量后验相对标准先验的KL项训练。评论家备份也从下一状态生成$m$个候选，用目标网络的
$\eta\min(Q_1^-,Q_2^-)+(1-\eta)\max(Q_1^-,Q_2^-)$
评分，其中$\eta\in[0,1]$越大，越偏重两个评论家中的较小值。对候选取最高评分，构造$r_t+\gamma$倍继续价值的回归目标；真终止时后续为0。扰动网络在限幅内提高$Q_1$，目标网络缓慢更新。部署实际输出式~\eqref{eq:rl-information-bcq}~选中的动作，不是行为模型。候选数$m$为正整数，增加它会增加生成与价值计算；增大$\Phi$则放宽偏离日志的范围\scite{rl4-fujimoto2019-bcq}

例如，标量教学状态中的生成器给出$-0.2,0.1,0.3$，扰动后为$-0.15,0.1,0.25$，对应$Q_1$评分$0.4,0.8,0.6$，于是执行0.1。这个筛选范围较小，却不说明三动作后果都已有充分记录：生成模型可能跨过数据空洞，扰动也可能越出支持。BCQ限制的是最大化搜索范围，覆盖条件仍须检查。

其他行为正则方法也限制偏离，但度量和施加位置不同。BEAR用最大均值差异比较行动者与行为动作样本，即比较核特征中的分布均值，约束备份尽量少查询分布外动作；距离含义取决于核、样本和阈值\scite{rl4-kumar2019-bear}
BRAC比较将行为分布散度加入策略目标或价值递推的不同方式。ReBRAC沿用TD3+BC，在行动者与评论家目标中分别调节行为惩罚，并检查网络、批量和归一化等实现因素。表~\ref{tab:rl-information-behavior-methods}~据此对照各自限制的位置；这些改动不增加真实信息\scite{rl4-wu2019-brac,rl4-tarasov2023-rebrac}

\begin{table*}[t]
  \centering
  \small
  \begin{tabularx}{\linewidth}{p{15mm}XXXX}
    \toprule
    方法 & 行为接近度 & 施加位置 & 执行动作来源 & 强度与代价 \\
    \midrule
    BEAR & 动作样本的核最大均值差异 & 行动者约束，间接限制备份 & 训练后的随机行动者；可对候选评分 & 阈值、核及拉格朗日权重；行为模型与样本比较 \\
    BRAC & 相对行为模型的KL等散度 & 策略正则，或同时进入价值递推 & 正则化行动者 & 散度权重；行为密度拟合误差 \\
    ReBRAC & 相对数据动作的平方距离 & 行动者及评论家目标的独立惩罚 & 确定性行动者 & 两处惩罚系数；数据下一动作与实现配置 \\
    \bottomrule
  \end{tabularx}
  \caption{行为正则的不同位置。表中“接近”都不表示已经获得未见动作的真实后果；实际比较还需固定动作尺度、数据与选模预算。}
  \label{tab:rl-information-behavior-methods}
\end{table*}

还可以不用显式行为距离，改用回报条件区分日志中的动作模式。Decision Transformer将过去状态、动作与剩余回报按因果顺序编码，以监督动作预测拟合$\pi_{\bm\phi}(a_t\mid h_t,g_t)$。这里又回到有限时域，其常见条件为未折扣剩余回报$g_t=\sum_{k=t}^{H-1}r_k$，不是默认折扣$G_t$。部署由用户指定希望达到的$g_0$，每步扣去实际奖励，按$g_{t+1}=g_t-r_t$更新。训练中$g_t$描述已发生轨迹的结果，部署中它是目标愿望；条件相同不保证结果可实现\scite{rl4-chen2021-dt}

\begin{example}[回报条件不能控制外部随机结果]
  \label{ex:rl-information-return-conditioning}
  设单步教学任务的奖励与动作独立：公平硬币以$1/2$概率给1，否则给0。仅取回报1的轨迹训练条件动作模型，可以完全拟合成功数据；但部署持续输入目标1，成功概率仍为$1/2$。按结果条件化只排除了失败样本，并未得到能够改变成功概率的动作。
\end{example}

回报条件模仿由结果标签区分行为，显式距离约束则惩罚分布或动作偏离，二者目标不同。前者可以表示高回报轨迹中的动作模式，却要面对幸运随机结果、不可达目标和日志外状态。判断执行策略是否改进，仍需第~\ref{sec:rl-information-evaluation}~节的价值证据。

% Outline: 4.4.3
\subsection{价值估计与策略更新中的外推控制}
\label{sec:rl-information-pessimism}

例~\ref{ex:rl-information-extrapolation}~中的未见动作被高估，除了限制策略，也可以改变价值训练。\term{保守Q学习}（conservative Q-learning, CQL）对缺少数据支持的高价值增加惩罚，使策略不容易因这些预测而偏离日志\scite{rl4-kumar2020-cql}
这里保持无限折扣价值和固定数据，采用有限动作的一种常用形式。从所用价值函数类选择$Q$，最小化Bellman平方残差与保守正则之和。正则先对全部有限动作价值取指数、求和再取对数，使较高值有较大影响；再减去数据动作的平均价值，约束数据外动作相对数据动作的乐观程度：
\begin{equation}
  \label{eq:rl-information-cql}
  \min_Q\
  \frac12\mathbb E_{\mathcal D}\bigl[Q(s,a)-y(s,a,r,s')\bigr]^2
  +\alpha\mathbb E_{s\sim\mathcal D}
       \left[\log\sum_{a'}e^{Q(s,a')}
                       -\mathbb E_{a\sim\widehat\mu_{\mathcal D}}Q(s,a)\right].
\end{equation}
其中$\alpha\geq0$控制正则强度，$y$由固定目标网络$Q^-$构造。评价当前$\pi^{(k)}$时，按其下一动作分布取期望，可令
$y=r+\gamma\mathbb E_{a'\sim\pi^{(k)}(\cdot\mid s')}Q^-(s',a')$；
控制备份也可以使用规定的动作最大化版本，真终止处都将继续价值置0。
$\widehat\mu_{\mathcal D}$是该状态的数据动作分布。为检查正则修正的方向，对动作价值求导，得到
$\alpha(\operatorname{softmax}(Q)(a\mid s)-\widehat\mu_{\mathcal D}(a\mid s))$；
于是无数据却预测很高的动作受到向下修正，数据动作得到相对支撑。连续动作不能直接枚举有限和，须按规定采样分布近似相应积分。

沿用例~\ref{ex:rl-information-extrapolation}~的动作与日志，另设教学预测
$Q(a_0)=0,Q(a_1)=4$。未见$a_1$的正则梯度约为$0.982\alpha$，梯度下降会压低它。这解释了局部抑制机制，却不证明训练后所有$Q(s,a)$都是可信下界。CQL理论还依赖表格或特定逼近、采样误差、正则强度及策略更新条件；常用版本控制策略期望价值下界，不能扩大为每个动作的逐点下界。

也可以改变继续价值的计算，避免价值训练查询新行动者提出的数据外动作。\term{隐式Q学习}（implicit Q-learning, IQL）从同一状态的数据动作价值中拟合偏向较高值的状态价值，再用它构造Q目标。要表达这种偏向，可用\term{期望分位数}（expectile）：以不对称权重计量高估和低估的平方误差，找出平均损失最小的位置。它不同于按概率质量切分的分位数，还考虑两侧偏差大小。对于平方可积的$X$和$\tau\in(0,1)$，比较所有实数位置$v$，定义
\begin{equation}
  \label{eq:rl-information-expectile}
  \ell_\tau(u)=|\tau-\mathbb I\{u<0\}|u^2,\qquad
  e_\tau(X)\in\argmin_v\mathbb E[\ell_\tau(X-v)].
\end{equation}
$\mathbb I$是指示函数，$u=X-v$为预测偏差。$\tau=1/2$使两侧权重相同，得到均值；$\tau>1/2$加重低估较大$X$的惩罚，使最优$v$向上移动。对离散分布求导，最优点满足
$\tau\mathbb E[(X-v)_+]=(1-\tau)\mathbb E[(v-X)_+]$。

IQL在每个状态按数据动作分布形成上述随机变量$Q(s,a)$。以下回到无限折扣价值，保持固定数据；$Q_{\bar{\bm\theta}}$为固定目标网络，$V_{\bm\psi}$为状态价值网络。先在各自参数类中最小化$L_V$拟合上侧位置，再最小化$L_Q$拟合奖励加继续价值，最后最小化$L_\pi$，用优势指数提高较好数据动作的拟合权重：
\begin{align}
  \label{eq:rl-information-iql}
  L_V(\bm\psi)
    &=\mathbb E_{(s,a)\sim\mathcal D}
       \ell_\tau(Q_{\bar{\bm\theta}}(s,a)-V_{\bm\psi}(s)),\nonumber\\
  L_Q(\bm\theta)
    &=\mathbb E_{\mathcal D}
       [Q_{\bm\theta}(s,a)-r-\gamma V_{\bm\psi}(s')]^2,\nonumber\\
  L_\pi(\bm\phi)
    &=-\mathbb E_{(s,a)\sim\mathcal D}
       [\exp(\beta\widehat A(s,a))\log\pi_{\bm\phi}(a\mid s)],\\
  \widehat A(s,a)&=Q_{\bar{\bm\theta}}(s,a)-V_{\bm\psi}(s).
  \nonumber
\end{align}
每次更新将目标量固定，真终止时$V(s')=0$。$\beta\geq0$是逆温度，控制优势差怎样转为相对拟合权重。可先完成价值拟合再提取策略，也可交替训练，但策略梯度不反向改变价值目标。实践中还使用双Q较小值和权重截断，须报告阈值。上侧expectile施加在同状态的动作价值上，而Q平方回归仍平均奖励及随机转移，不挑选幸运反馈。原IQL用分开的$V$、$Q$目标保持了可选择动作与不可选择环境结果的区别\scite{rl4-kostrikov2022-iql}

\begin{example}[三个动作的IQL更新与提取]
  \label{ex:rl-information-iql}
  在教学状态$s$，三个数据动作等频出现，目标Q值为0、1、3，取$\tau=3/4$。为找状态价值，先假设$v\in[1,3]$，将一项低估与两项高估按损失权重平衡：
  \[
    \tfrac34(3-v)=\tfrac14[(v-0)+(v-1)],
  \]
  解为$v=2$，确实落在假设区间；损失凸，故为全局最优点。对到达$s$的非终止转移，若奖励0.5、$\gamma=0.9$，继续价值便给出Q回归目标$0.5+0.9(2)=2.3$。

  提取策略时取$\beta=\log2$。三个优势$-2,-1,1$对应未归一化权重$1/4,1/2,2$。策略若可在所有分类分布中选择，则最大化加权对数似然等价于按等频数据的权重归一化，得到
  \[
    \pi(a_1\mid s)=1/11,\quad
    \pi(a_2\mid s)=2/11,\quad
    \pi(a_3\mid s)=8/11.
  \]
  若这些Q恰好等于单步真实奖励，提取策略的真实价值为$26/11$，高于等频行为的$4/3$，仍低于总选$a_3$的3。这一特殊情形也说明，拟合所得$V=2$不等于执行策略价值$26/11$。
\end{example}

\begin{algorithm}[固定经验上的IQL与策略提取]
  \label{alg:rl-information-iql}
  \begin{algorithmic}[1]
    \Require 固定$\mathcal D$，$\gamma,\tau,\beta$，价值与策略更新预算，权重截断阈值
    \Ensure 可执行$\pi_{\bm\phi}$及送往独立评价的数据职责说明
    \State 初始化$Q_{\bm\theta},Q_{\bar{\bm\theta}},V_{\bm\psi},\pi_{\bm\phi}$
    \For{每次价值更新}
      \State 从$\mathcal D$抽取转移批次，按$L_V$拟合数据动作Q值的expectile
      \State 按真终止标记屏蔽继续价值，构造固定目标$r+\gamma V_{\bm\psi}(s')$
      \State 按$L_Q$更新评论家，再缓慢更新目标Q参数
    \EndFor
    \For{每次策略提取更新}
      \State 固定价值网络，计算数据动作优势和截断后的指数权重
      \State 最小化加权负对数似然$L_\pi$
    \EndFor
    \State 固定提取策略，按第~\ref{sec:rl-information-evaluation}~节评价，不以$V$直接充当最终证据
  \end{algorithmic}
\end{algorithm}

算法~\ref{alg:rl-information-iql}~按给定更新预算先拟合价值，再用固定价值提取动作分布，最后将执行策略送往评价。表~\ref{tab:rl-information-extraction}~进一步对照提取方式：同一价值函数经不同动作来源形成不同策略，其逼近误差与覆盖风险也不同。

\begin{table*}[htbp]
  \centering
  \small
  \begin{tabularx}{\linewidth}{p{26mm}XXX}
    \toprule
    提取方式 & 动作来源与价值查询 & 逼近误差 & 覆盖风险 \\
    \midrule
    直接贪心 & 全动作优化或行动者提出动作；查询其Q值 & 动作优化不精确、评论家误差 & 容易主动寻找数据外高估动作 \\
    优势加权拟合 & 用数据动作训练执行分布；训练查询数据Q值 & 策略族与加权分布不匹配、优化误差 & 连续输出仍可能落在稀疏间隙 \\
    行为采样后筛选 & 生成$m$个候选，按Q取最大或重采样 & 行为生成偏差、有限$m$误差 & 生成概率非零不证明环境后果有支持 \\
    \bottomrule
  \end{tabularx}
  \caption{价值学习之后的三个动作接口。增大候选数或策略容量可能减少提取误差，也可能更充分地利用价值外推误差。}
  \label{tab:rl-information-extraction}
\end{table*}

扩散模型也可接入这些动作提取方式。Diffusion-QL拟合条件动作扩散模型，并在去噪目标外最大化生成动作的价值，让价值梯度经生成过程改变策略。行为建模与价值驱动训练由此结合，去噪步数则同时影响训练和推理成本\scite{rl4-wang2023-diffusionql}
IDQL研究价值目标对应的行为重加权分布，以扩散行为模型生成候选，按评论家导出的权重重采样，或采用贪心提取变体。它试图减少简单行动者无法表示多模态隐含策略的误差；生成候选仍不提供新的环境后果证据\scite{rl4-hansenestruch2023-idql}

图~\ref{fig:rl-information-extraction}~区分价值梯度改变生成器与价值筛选已有候选。流模型也能生成候选或给出可微动作映射；要据此说明改进，仍须接入有依据的价值加权、策略目标或规划准则。单独拟合条件生成模型没有完成这一比较步骤。

\begin{figure}[htbp]
  \centering
  % Two distinct value interfaces for conditional action generation.
\begin{tikzpicture}[
  x=\linewidth,y=1mm,
  font=\figurefont\fontsize{8}{10}\selectfont,
  box/.style={draw=BookRule,fill=BookPaper,line width=.8pt,
    text width=.23\linewidth,minimum height=10mm,align=center,inner sep=1mm},
  flow/.style={-{Latex[length=1.8mm]},draw=BookInk,line width=.9pt},
  gradient/.style={-{Latex[length=1.8mm]},draw=BookAmber,line width=.8pt,dashed}
]
  \node[anchor=west] at (.01,12) {(a) 价值驱动生成器训练};
  \node[box,fill=BookTeal!8] (generator) at (.14,0) {条件生成器};
  \node[box] (candidate) at (.5,0) {生成动作$a$};
  \node[box] (value) at (.86,0) {价值$Q(s,a)$};
  \draw[flow] (generator) -- (candidate);
  \draw[flow] (candidate) -- (value);
  \draw[gradient] (value.south) -- (.86,-14)
    -- node[below] {价值梯度经生成过程回传} (.14,-14) -- (generator.south);
  \node[anchor=west] at (.01,-30) {(b) 行为生成后的价值筛选};
  \node[box,fill=BookTeal!8] (behavior) at (.14,-43) {行为生成器};
  \node[box] (candidates) at (.5,-43) {$a_1,\ldots,a_m$};
  \node[box,fill=BookBlue!8] (selection) at (.86,-43) {加权重采样\\或取最大};
  \node[box] (critic) at (.5,-66) {价值评分};
  \node[box] (action) at (.86,-66) {最终动作};
  \draw[flow] (behavior) -- (candidates);
  \draw[flow] (candidates) -- (selection);
  \draw[flow] (candidates) -- (critic);
  \draw[flow] (critic.east) -- (.7,-66) -- (.7,-55) -- (selection.south west);
  \draw[flow] (selection) -- (action);
\end{tikzpicture}

  \caption{生成式动作的两种价值接口。实线为部署时的候选与动作流，虚线为训练中从价值目标进入生成器的梯度；两条路径不应混称为同一种“扩散策略改进”。}
  \label{fig:rl-information-extraction}
\end{figure}

动作还可以组合为序列。近期DEAS通过与显式行动者分离的价值学习，减轻动作块评价的过估计，其ICLR 2026正式版本研究长时程离线任务与动作序列模型。动作块减少显式选择次数，却扩大了组合动作空间；固定长度序列也不等于能够中途依反馈行动的Option，更不说明日志覆盖全部动作块。这里只将DEAS作为该接口的近期研究，基准表现不构成普遍最优性结论\scite{rl4-kim2026-deas}

% Outline: 4.4.4
\subsection{学习模型上的规划与策略优化}
\label{sec:rl-information-model-learning}

第~\ref{chap:rl-learning}~章允许新交互，因此规划器查询到模型的可疑预测后，还可采样校正；固定经验下无法这样验证。我们仍可从日志拟合$\widehat R(s,a)$和$\widehat P(s'\mid s,a)$，并计算不确定性分数$u(s,a)$。留出数据能检查其覆盖区域中的部分预测误差；数据外的真实后果却不能仅由小误差、小集成分歧或高似然证明。

MOPO在模型滚动时扣除不确定性惩罚：从日志状态出发，按当前策略在$\widehat P$中生成有限长度轨迹，以
$\widetilde R=\widehat R-\lambda u$作为模拟奖励，再混合真实与模拟转移更新策略并重复这一过程。模型只由真实$\mathcal D$拟合，模拟转移保留来源标识。滚动长度控制误差可累积的范围，$\lambda$调节不确定区域的惩罚\scite{rl4-yu2020-mopo}

要把惩罚后的回报解释为真实价值下界，需要知道它扣除的量能否覆盖模型误差。固定$\pi$和无限折扣目标，令$\widehat d^\pi$为$\widehat P$下的归一化占用。为计算真实价值减去模型价值，将两个Bellman方程相减，并按模型轨迹的访问概率累积，得到
\begin{equation}
  \label{eq:rl-information-model-telescoping}
  J_M(\pi)-J_{\widehat M}(\pi)
  =\frac1{1-\gamma}\mathbb E_{\widehat d^\pi}
       \left[R-\widehat R
                  +\gamma(P-\widehat P)V_M^\pi\right].
\end{equation}
这里$(P-\widehat P)V$是两核给出的下一价值期望之差。共同$\rho_0$和有界价值使折扣尾项消失；占用已归一化，故保留$1/(1-\gamma)$。若在该策略访问区域内，
$|R-\widehat R|\leq u_R$且
$|(P-\widehat P)V_M^\pi|\leq u_P$，
则每处误差项不低于$-u_R-\gamma u_P$。按同一非负占用分布平均，得到扣除奖励与转移价值误差上界后的保守回报：
\[
  J_M(\pi)\geq
    \frac1{1-\gamma}\mathbb E_{\widehat d^\pi}
       [\widehat R-u_R-\gamma u_P].
\]
因此，$\lambda u$只有与实际误差上界建立联系，才具有上述下界含义。用它选择策略时，条件还须同时覆盖参与选择的候选，不能沿用单个预固定策略的界。所有集成模型也可能一致地学错；没有覆盖或正确结构支撑时，模型分歧仍只是启发式诊断。

除惩罚奖励外，MOReL在不可信状态动作处转入低回报吸收状态，使规划器少离开可信区域\scite{rl4-kidambi2020-morel}
COMBO对模型滚动状态动作施加保守价值正则，不必先构造同一种显式不确定性分数\scite{rl4-yu2021-combo}
两者都仍需要真实数据和各自保证中的模型误差条件。表~\ref{tab:rl-information-model-methods}~在相同日志缺口下比较这些选择：未见动作若实际很好，保守限制会放弃可实现回报；若实际很差，限制可能阻止规划器利用错误预测。

\begin{table*}[t]
  \centering
  \small
  \begin{tabularx}{\linewidth}{p{17mm}XXXX}
    \toprule
    方法 & 限制位置 & 候选轨迹来源 & 真实数据依赖 & 保守程度与代价 \\
    \midrule
    MOPO & 模型奖励减去误差代理惩罚 & 日志状态出发的短模型滚动 & 动力学拟合及误差校准 & 惩罚系数、滚动长度；可能拒绝高回报稀疏区域 \\
    MOReL & 不可信区域进入低回报吸收态 & 悲观模型中的规划或学习 & 可信区域判定与模型置信 & 阈值与吸收惩罚；错误拒绝会缩小可达策略类 \\
    COMBO & 真实与模型数据上的保守价值正则 & 学习模型中的策略滚动 & 模型拟合及数据动作的价值支撑 & 正则及混合比例；过强会压低真实可行动作 \\
    \bottomrule
  \end{tabularx}
  \caption{模型法限制外推的三个位置。比较成立的前提是任务、日志、模型容量与查询预算相同；减少显式不确定性估计不等于不再需要误差依据。}
  \label{tab:rl-information-model-methods}
\end{table*}

另一种模型使用方式是直接生成整段轨迹。Trajectory Transformer自回归地预测状态、动作与奖励序列，规划时用束搜索保留高评分候选，再执行选中轨迹的首个动作；收到真实观测后重新搜索。改进还取决于搜索评分与候选可行性，不能仅由下一token预测损失推出。
Diffuser按当前状态等条件，通过回报引导的扩散采样形成计划，也可执行首个动作后滚动重规划。前者搜索自回归候选，后者引导去噪生成计划，但两者轨迹都是模型预测，没有增加观测经验\scite{rl4-janner2021-trajectory,rl4-janner2022-diffuser}

部署接收反馈并重规划，不说明训练期间有在线调参权限。若为选超参数反复在真实环境执行计划，这些查询须计入第~\ref{sec:rl-information-limited-interaction}~节的新增预算。搜索与模型滚动的通用计算见第~\ref{chap:rl-planning}、\ref{chap:rl-learning}~章；固定日志下，规划器仍不能自行验证数据外预测。

% Outline: 4.4.5
\subsection{策略改进的保证与比较}
\label{sec:rl-information-improvement}

前述三类限制都产生了候选，但候选价值超过估计基线还不等于真实改进。需要先规定比较哪些策略，以及估计误差能否同时约束它们。下面固定候选集合与基线，用共同误差事件计算选择后策略最多可能损失多少回报，不依赖特定网络或优化器。

\begin{theorem}[共同误差事件上的基线性能控制]
  \label{thm:rl-information-improvement}
  固定任务、初始分布、时域口径、候选策略集合$\Pi$及基线$\pi_b$。在共同事件$\mathcal E$上，假设每个
  $\pi\in\Pi\cup\{\pi_b\}$均有
  $|\widehat J(\pi)-J(\pi)|\leq\varepsilon$。从$\Pi$中选出的策略$\widehat\pi$若满足
  $\widehat J(\widehat\pi)\geq\widehat J(\pi_b)$，则该事件中真实价值满足
  \[
    J(\widehat\pi)\geq J(\pi_b)-2\varepsilon.
  \]
  若$\mathbb P(\mathcal E)\geq1-\delta$，结论也以至少$1-\delta$的概率成立。
\end{theorem}

\begin{proof}
  为把估计上的选择规则转为真实价值比较，在共同事件内插入候选和基线的估计值：
  \[
  \begin{aligned}
    J(\widehat\pi)
    &\geq\widehat J(\widehat\pi)-\varepsilon\\
    &\geq\widehat J(\pi_b)-\varepsilon\\
    &\geq J(\pi_b)-2\varepsilon.
  \end{aligned}
  \]
  第一、三步分别使用$\widehat\pi$与$\pi_b$的共同误差界，第二步使用选择规则。事件$\mathcal E$发生时这条链必成立，故结论概率不低于$\mathbb P(\mathcal E)$。
\end{proof}

这一界允许真实回报比基线低$2\varepsilon$，并不保证严格提高。将选择阈值提高到
$\widehat J(\widehat\pi)-\widehat J(\pi_b)>2\varepsilon$，同样推导才给出严格正的真实增量。共同事件是关键条件：预固定单策略的95\%区间不能直接约束从许多候选中选出的最好者；候选集若由评价数据生成，还须处理适应性。

也可以在候选形成时保留基线对稀少数据区域的行为。SPIBB用基线自举实现这一约束，即数据少处沿用基线，其他动作才重新分配概率。有限MDP中以计数阈值$N_{\min}$识别稀少对，令
$\mathcal B=\{(s,a):N(s,a)<N_{\min}\}$，
其中$N(s,a)$为数据计数。要求$\mathcal B$内$\pi(a\mid s)=\pi_b(a\mid s)$，在其余动作上优化经验模型价值。例如基线动作概率$(0.6,0.3,0.1)$的第三项稀少，则保持0.1，而不是删去它再归一化；剩下0.9的概率可以在前两项间调整。

原始$\Pi_b$-SPIBB界依赖有限状态动作、有界奖励、已知且可执行的基线、按计数建立的模型置信条件及规定的受约束求解，控制的是相对基线容许损失。损失还随置信半径、折扣累积和经验模型增益变化，所以名称中的“安全”不表示每次运行超过基线。神经网络近似、模型误设或基线概率错误都不能直接继承表格结论\scite{rl4-laroche2019-spibb}

\begin{table*}[t]
  \centering
  \small
  \begin{tabularx}{\linewidth}{p{28mm}XXX}
    \toprule
    共同报告项 & 行为克隆 & 可复现行为策略 & 等预算离线改进方法 \\
    \midrule
    价值证据 & 留出OPE或独立测试；不是拟合准确率 & 原始版本、真实测试或日志回报及其条件 & 同一评价器、置信处理和查询权限 \\
    偏离与悲观设置 & 模型容量、拟合与动作输出形式 & 原收集分布与版本 & 行为惩罚、价值惩罚、模型可信阈值 \\
    训练成本 & 数据遍历、批量、更新次数 & 复现基线所需成本单列 & 评论家、生成器、模型及调参总成本 \\
    推理成本 & 一次动作预测或采样 & 原控制器调用成本 & 候选数、去噪步数、搜索与滚动长度 \\
    \bottomrule
  \end{tabularx}
  \caption{方法比较的报告口径，而非虚构的实验结果表。真实测试、离线估计和理论性能界分别标注，不能放在同一列中当作等强度证据。}
  \label{tab:rl-information-reporting}
\end{table*}

表~\ref{tab:rl-information-reporting}~将训练收益、证据类型及成本分别报告。基线本身可能回报很低，也可能违反业务安全要求，因此接近基线既不表示接近全局最优，也不说明满足碰撞概率或资源成本约束。本节只控制给定条件下的相对回报；执行安全和约束保证由第~\ref{chap:rl-trustworthy}~章展开。

% Outline: 4.5
\section{有限新增经验下的策略修正}
\label{sec:rl-information-limited-interaction}

回到定理~\ref{thm:rl-information-unseen-action}~，原日志只含$a_0$的零奖励，无法区分两环境。若权限放宽，执行一次$a_1$，确定性奖励$+1$或$-1$便能排除一个解释。这里新增的是实际动作后果，而非一次模型预测。奖励有噪声时，一次查询通常还不能辨清环境，获得查询权限也不意味着估计误差已经可忽略。

这种信息补充必须服从每轮协议：开始前声明可执行策略、步数、重置次数、专家帮助权限和验证用途。专家动作与主动随机探索的选择机制不同，新增日志要记录实际机制。即使用于最终报告而不参与梯度更新，试跑也消耗真实交互预算。

新经验加入后，原离线保守构造也要随证据重查。行为约束过强会挡住已有新数据支持的动作，突然取消又可能让尚未校准的评论家驱动危险外推。AWAC将离线与在线经验用于优势加权动作拟合，减少对显式行为密度的依赖\scite{rl4-nair2021-awac}
IQL也可将新经验接入原有$V,Q$与加权策略计算。Cal-QL则处理保守初始化的数值尺度：在下估当前策略价值时，用参考策略价值限制过度压低，减少微调初期仅为恢复尺度所作的更新。这是价值尺度校准，不是任意任务的概率校准保证\scite{rl4-nakamoto2023-calql}

\begin{table*}[t]
  \centering
  \small
  \begin{tabularx}{\linewidth}{p{18mm}XXXX}
    \toprule
    起点 & 新旧经验混合 & 评论家变化 & 行为限制变化 & 探索恢复要点 \\
    \midrule
    AWAC & 旧日志与在线回放用于动态规划和优势拟合 & 用当前继续策略的价值更新优势 & 加权似然的隐式接近度随数据改变 & 行动者必须真正产生可补足覆盖的新动作 \\
    IQL微调 & 新转移加入价值和策略训练 & 数据动作expectile与Q目标随新分布变化 & 温度、截断和混合比例决定偏离程度 & 不查询数据外Q不等于主动探索已经充分 \\
    Cal-QL & 离线初始化后混合在线经验 & 参考价值校准限制过度低估 & 保守强度与新证据一起检查 & 减轻初始价值尺度失配，仍需覆盖新增行为 \\
    \bottomrule
  \end{tabularx}
  \caption{有限交互修正中的组件职责。三行不是新的互斥算法分类；混合比例、约束强度和探索预算都需报告，不能只比较最终训练步数。}
  \label{tab:rl-information-finetuning}
\end{table*}

表~\ref{tab:rl-information-finetuning}~列出的组件可以接入同一轮修正。算法~\ref{alg:rl-information-round}~先收集并更新数据，再形成候选，最后重新评价。新经验同时改变覆盖与行为分布，所以旧评价器的概率和误差条件也需要重新核对。

\begin{algorithm}[一轮有限经验修正]
  \label{alg:rl-information-round}
  \begin{algorithmic}[1]
    \Require 当前日志$\mathcal D^{(k)}$、基线、候选，收集预算$B_k$与独立评价预算$E_k$
    \Ensure 更新后的候选集合与本轮可支持的性能结论
    \State 固定本轮收集策略及版本，声明可访问状态、重置和帮助权限
    \State 在$B_k$内收集经验，保存策略概率、轨迹边界及用途标识
    \State 构成$\mathcal D^{(k+1)}$，重新检查覆盖及新旧数据混合权重
    \State 更新表示、价值、行为或环境模型，按选定接口提取候选
    \State 固定待比较策略，在保留日志或$E_k$内的独立交互上重新评价
    \State 依据置信规则保留、替换或拒绝候选，记录下一轮需要补充的信息
    \State 扣除收集、验证、重置及专家调用预算；预算用尽时停止
  \end{algorithmic}
\end{algorithm}

除了总预算，协议还决定动作能否根据上一批反馈改变。下面在相同两次查询下，只改变反馈可以进入决策的轮次。

\begin{example}[相同两次查询与不同适应轮次]
  \label{ex:rl-information-batches}
  沿用两个单步环境：$a_0$奖励0，$a_1$分别恒为$+1$或$-1$。为教学比较，假设先验各为$1/2$，共允许两个新回合，每回合一次动作。

  两次动作若都须在取得反馈前确定，预先选几次$a_1$都使两环境的平均收集回报为0。允许分两轮时，可先执行$a_1$，正奖励后再选$a_1$，负奖励后改选$a_0$；两环境总回报为2和$-1$，先验平均为$1/2$。
\end{example}

这项差距属于收集期间总回报，不是所有最终指标的差距。若只要求收集结束后确定动作，一次性设计中安排一次$a_1$也可辨认环境。相同样本预算配合不同反馈时机，仍可能产生不同的行动过程，因此多轮方法须报告原日志、每轮新增预算和反馈使用时机，并将反复测试的选择偏差纳入证据处理。

% Outline: 4.6
\section{两类信息限制叠加时的强化学习}
\label{sec:rl-information-combined}

% Outline: 4.6.1
\subsection{两类限制下的决策目标与可用信息}
\label{sec:rl-information-combined-target}

此前讨论了怎样用历史表示处理执行时状态不可见，离线评价则暂时假设日志记录完整状态。现在撤掉完整状态假设并保留固定经验：日志中的历史未必包括收集者实际使用的所有线索。因此，不能只将评价式中的$s_t$换成$\bm z_t$。需要先固定目标量，再分别核对执行策略可用的信息、日志保存的信息和行为策略当时使用的信息。

以下恢复有限时域$H$，固定$\rho_0$、$0\leq\gamma<1$和执行可得历史$h_t$，以$\Pi$表示允许的历史策略类。可以针对一个$J_H(\pi)$、两个候选的差或相对基线增益提出问题；若要找$\Pi$内的最优行为，就需要比较全部相关候选。一个候选可评价，只回答该候选的问题，不足以推出能够找到全局最优。

\begin{table*}[htbp]
  \centering
  \small
  \begin{tabularx}{\linewidth}{XXX}
    \toprule
    执行可用信息 & 日志实际字段 & 行为策略曾使用的信息 \\
    \midrule
    只见位置与历史动作、奖励 & 同样的全回合历史与边界 & 同样历史及独立随机数：可建立历史条件校正 \\
    可以保留起始提示 & 日志只保留最近窗口 & 收集者记住提示：执行输入与记录信息不匹配 \\
    只见传感观测 & 记录动作概率但缺少概率所条件化的上下文 & 专家还看到未记录状态：不能把边缘频率当成随机分配机制 \\
    回合开始重置记忆 & 片段无回合标识、无前缀 & 行为控制器跨片段保留记忆：初始化机制未知 \\
    \bottomrule
  \end{tabularx}
  \caption{离线POMDP中的三套信息。动作概率需要连同其条件信息和收集协议解释；一个额外字段不能自动修复全部缺口。}
  \label{tab:rl-information-information-sets}
\end{table*}

表~\ref{tab:rl-information-information-sets}~第二行的日志缺少部署可保留的提示，因而无法构造待评价策略的输入；第三行缺少专家选动作的依据，还可能改变动作分组与结果之间的统计关系。这两种遗漏都涉及不可见信息，却阻断评价的不同步骤。要判断模型是否足够，应按当前价值或动作比较检查，而非要求恢复所有隐藏状态名称。

\begin{definition}[给定问题上的决策等价]
  \label{def:rl-information-decision-equivalence}
  决策等价以给定的允许策略类、时域、初始条件与目标结论为条件。两个潜在模型若对该问题涉及的全部价值或动作比较给出相同结论，则在该问题上决策等价。该等价不要求隐藏状态名称一致，也不延伸到未声明的新奖励或策略类。
\end{definition}

例如，交换隐藏状态名称会改变参数表示，却不改变任何可执行策略的价值，控制无需唯一的隐藏语义。两个同样预测原日志的模型则可能在强制执行别的动作后给出不同后果；这种差异会改变策略判断。先声明目标结论，才能确定哪些模型差异必须识别。

% Outline: 4.6.2
\subsection{策略结论的可识别程度}
\label{sec:rl-information-identification-levels}

以表~\ref{tab:rl-information-information-sets}~为起点，离线POMDP至少要核对三个条件。历史应保留执行输入并具有正确回合边界；行为应只依据记录历史和独立随机化，或满足另一个明确识别条件；目标历史动作应有行为支持。物理状态隐藏不等于行为机制不明：控制器若也只见记录历史，隐藏状态虽影响反馈，却没有额外进入动作选择。

日志未必能对每个目标都满足这些条件。为了说明已有信息到底排除了哪些解释，可以把具有同一可观测总体分布、又满足既定结构假设的模型保留下来，再查看它们对目标量给出了哪些答案。

\begin{definition}[一致模型集合与识别集合]
  \label{def:rl-information-identified-set}
  给定可观测日志总体分布$p_{\mathrm{obs}}$与结构假设集合$\mathcal A_0$，一致模型集合定义为
  \[
    \mathcal M(p_{\mathrm{obs}},\mathcal A_0)
       =\{M:M\text{满足}\mathcal A_0
                  \text{且产生}p_{\mathrm{obs}}\}.
  \]
  每个$M$同时指定潜在环境与相容的数据生成机制。将目标量$T(M)$在这些模型中的取值收集为识别集合
  $\mathcal I_T=\{T(M):M\in\mathcal M\}$，假设$\mathcal M$非空。
  $\mathcal I_T$为单元素时，$T$被点识别；若非单元素但严格包含在预先允许范围内，则得到部分识别。对于策略排序，需单独检查所有允许模型是否给出相同次序；若存在相反排序，则排序不可识别。
\end{definition}

确定数值与确定决策是不同要求。例如，性能差$T(M)=J_M(\pi_1)-J_M(\pi_2)$若原可取$[-10,10]$，日志和结构将其缩小到$[-1,2]$，便已提供部分数值信息，但正负值仍并存，不能排序。若缩小到$[1,2]$，数值仍未唯一确定，符号却足以说明$\pi_1$较好。下面按目标是否点识别、是否只有部分范围、是否仍不能回答决策展开；部分识别本身不等于无法决策。

一致模型和识别集合描述总体分布已知时的限制，有限日志还会引入采样误差，需要再构造相应目标量或识别集合的统计置信集合。一个模型拟合有限日志很好，既没有排除其他总体相容模型，也没有单独控制采样误差。

行为机制是否含有额外隐藏依据，是其中一个关键区别。\term{隐藏混杂}（hidden confounding）指未记录变量同时影响动作选择与后果，观察到的动作组因而不能直接代表随机试验组。图~\ref{fig:rl-information-confounding}~对照普通隐藏状态与这种共同原因；箭头是声明的因果生成假设，不是拟合日志自动发现的事实。

\begin{figure*}[htbp]
  \centering
  % Causal assumptions, not dependencies inferred from the log.
\begin{tikzpicture}[
  x=\linewidth,y=1mm,
  font=\figurefont\fontsize{8}{10}\selectfont,
  observed/.style={draw=BookBlue,fill=BookBlue!6,line width=.8pt,
    minimum width=12mm,minimum height=9mm,align=center},
  hidden/.style={observed,draw=BookMuted,fill=white,dashed},
  cause/.style={-{Latex[length=1.8mm]},draw=BookInk,line width=.9pt}
]
  \node[anchor=west] at (.015,32) {(a) 普通隐藏状态};
  \node[anchor=west] at (.56,32) {(b) 隐藏混杂};
  \node[hidden] (u1) at (.23,17) {$U$};
  \node[observed] (h1) at (.06,0) {$h_t$};
  \node[observed] (a1) at (.23,0) {$a_t$};
  \node[observed] (y1) at (.4,0) {$Y$};
  \draw[cause] (u1) -- (h1);
  \draw[cause] (u1) -- (y1);
  \draw[cause] (h1) -- (a1);
  \draw[cause] (a1) -- (y1);
  \node[hidden] (u2) at (.77,17) {$U$};
  \node[observed] (h2) at (.6,0) {$h_t$};
  \node[observed] (a2) at (.77,0) {$a_t$};
  \node[observed] (y2) at (.94,0) {$Y$};
  \draw[cause] (u2) -- (h2);
  \draw[cause] (u2) -- (y2);
  \draw[cause] (h2) -- (a2);
  \draw[cause] (a2) -- (y2);
  \draw[cause,BookAmber,line width=1.2pt] (u2) -- (a2);
  \node[align=center,text=BookMuted] at (.23,-17)
    {给定历史后\\行为随机数独立于$U$};
  \node[align=center,text=BookMuted] at (.77,-17)
    {给定历史后\\$U$仍影响动作};
\end{tikzpicture}

  \caption{隐藏状态与隐藏混杂的区别。两图中的$U$均未记录，$h_t$为记录历史，$Y$为未来反馈。左图假设行为仅由$h_t$和独立随机数决定；右图多出$U\to a_t$，给定$h_t$仍不能阻断动作与后果之间的未记录共同原因。箭头表示因果假设。}
  \label{fig:rl-information-confounding}
\end{figure*}

\begin{example}[同样日志与相反的动作价值]
  \label{ex:rl-information-confounding}
  构造单步教学任务，观测$o$恒定，类型$U\in\{0,1\}$等概率且未记录，专家却能看到它。下面给出两个模型中专家的动作和每种类型执行两动作的潜在奖励：
  \[
  \begin{array}{cc|c|rr}
    \text{模型}&U&\text{专家动作}&r(a_0)&r(a_1)\\
    \hline
    M_+&0&a_0&1&3\\
    M_+&1&a_1&-1&1\\
    M_-&0&a_1&3&1\\
    M_-&1&a_0&1&-1
  \end{array}
  \]
  两模型都以$1/2$记录$(o,a_0,1)$、以$1/2$记录$(o,a_1,1)$。部署看不到$U$，若始终选一个固定动作，其价值须按总体类型概率平均该动作后果，而非只平均专家选它的类型：
  \[
  \begin{array}{c|rr}
    &J_1(\pi_{a_0})&J_1(\pi_{a_1})\\
    \hline
    M_+&(1-1)/2=0&(3+1)/2=2\\
    M_-&(3+1)/2=2&(1-1)/2=0
  \end{array}
  \]
  因此，相同记录分布在两模型中对应相反的最优固定动作。
\end{example}

只用$o$拟合动作频率，会得到$\widehat\mu(a_0\mid o)=\widehat\mu(a_1\mid o)=1/2$。用它作普通IS，对任一固定动作目标都算得
$(1/2)\times2\times1=1$，既不是0也不是2。边缘上两动作都出现，并未恢复条件随机化：专家选该动作的样本只来自特定类型，不能代表总体类型执行该动作的结果。

有时可以用代理变量补充对隐藏量的间接观测，或增加传感数据，但识别仍需明确的条件独立、可逆性或完备性等结构。字段变多不说明这些条件成立。部分POMDP离策略评价工作利用代理观测和矩阵秩条件处理隐藏机制，不能将这类结论改写为训练循环网络即可消除混杂\scite{rl4-tennenholtz2020-pomdp-ope}

% Outline: 4.6.3
\subsection{完全可识别时的历史表示与决策学习}
\label{sec:rl-information-full-identification}

先考虑不含上述隐藏混杂的一组条件：日志从回合开始完整保留$h_t$，行为以已知$\mu_t(a_t\mid h_t)$随机选择，给定历史后的随机数与潜在环境独立；目标也只使用记录历史，初始分布和环境共同，目标轨迹有行为支持。于是每个历史层内的动作按已知机制分配，物理状态隐藏不再给行为增加未知依据。我们可以据此计算目标与行为的观测轨迹比，获得点识别目标价值所需的校正。

为将第~\ref{sec:rl-information-estimators}~节的校正改到历史输入，写观测轨迹为
$\tau^o=(o_0,a_0,r_0,\ldots,o_H)$，
将给定历史和动作后的反馈分布记作$q(y_{t+1}\mid h_t,a_t)$。上述行为条件使目标与行为共享这个反馈因子；按时间分解密度，以目标除以行为，共同初始和反馈项约去：
\begin{equation}
  \label{eq:rl-information-history-is}
  \frac{p_\pi(\tau^o)}{p_\mu(\tau^o)}
   =\frac{p(o_0)\prod_{t=0}^{H-1}
                  \pi_t(a_t\mid h_t)q(y_{t+1}\mid h_t,a_t)}
          {p(o_0)\prod_{t=0}^{H-1}
                  \mu_t(a_t\mid h_t)q(y_{t+1}\mid h_t,a_t)}
   =\prod_{t=0}^{H-1}\frac{\pi_t(a_t\mid h_t)}
                              {\mu_t(a_t\mid h_t)}.
\end{equation}
隐藏状态已在共同反馈因子中积分，因此最终只保留历史条件动作比。也可以先在完整隐藏轨迹上消去共同环境项，再求和隐藏状态；由于两策略的动作概率都只依赖可见历史，才能将它们提出该求和。

这一消去不能通过将$h_t$随意替换为$o_t$保留。行为若记住早期提示，只按当前画面拟合出的边缘频率就不是实际动作条件概率。专家若还看到例~\ref{ex:rl-information-confounding}~中的$U$，其动作项会随隐藏类型变化，无法从隐藏状态求和中提出，校正便失去依据。

识别条件成立后，可以直接按完整历史评价，或用有依据的充分表示降低计算成本。正确POMDP已知时用精确信念，预测结构可识别时可用PSR；潜在模型、循环或注意力记忆则需要检查允许动作上的预测或控制充分性。相应地，FQE可以在有限时域历史树上回归，行为约束作用于$\pi(a\mid h)$，提取动作时查询同一历史或经过论证的充分表示。直接保留历史减少了压缩假设，却可能使历史动作组合非常稀疏。

这些计算仍有表示构造、行为概率估计、有限轨迹采样与优化误差，报告时应分开。训练用burn-in恢复记忆，部署却每$L$步清零，会变成不同策略；训练读到回合开始提示，部署只有窗口，也不是原来评价的对象。总体可识别排除了信息上的不可能性，但有限数据、所选网络和优化还未必实现相应估计。

早期提示任务可以将表示错误与不可识别分开。假设行为在每种提示下随机尝试两个出口，完整历史也被记录，则动作价值在总体上可识别。一个编码器保留$c$，另一个只重建走廊图像；延迟很长时，提示帧占平均重建损失的比例很小，两者重建误差可能相近。按原例不折扣成功奖励，前者仍可达到1，后者至多$1/2$。此时失败来自编码遗漏，不能解释为环境本来无法识别。

% Outline: 4.6.4
\subsection{部分可识别时的历史表示与决策学习}
\label{sec:rl-information-partial-identification}

若同一日志和结构仍容许多个价值，就需要保留定义~\ref{def:rl-information-identified-set}~中的一致模型集合$\mathcal M$，而不是选一个拟合模型当作唯一真相。对于固定$\pi$，比较$\mathcal M$内各模型的价值，取下确界和上确界，计算
$[\inf_MJ_M(\pi),\sup_MJ_M(\pi)]$，
作为识别集合的上下包络；也可改变允许的混杂强度，检查敏感性范围。包络只描述最小和最大边界，未必每个中间值都可实现。有限日志还须在外层加入总体分布未知的采样误差。

排序所需的信息可能比单独价值少。固定两个候选，在每个相同模型中先计算$\pi_1$价值减去$\pi_2$价值，再对$\mathcal M$取下确界。若这个最小边界仍为正，全部允许模型便给出同一排序；只看平均模型不够，因为另一个同样符合日志的模型仍可能反转次序。

\begin{theorem}[共同模型下的性能差确定排序]
  \label{thm:rl-information-partial-order}
  给定非空一致模型集合$\mathcal M$，假设两个策略在每个模型中的价值均存在。若
  \[
    \inf_{M\in\mathcal M}
          [J_M(\pi_1)-J_M(\pi_2)]>0,
  \]
  则所有允许模型都给出$\pi_1$优于$\pi_2$的排序，两个单独价值区间即使重叠也不影响结论。
\end{theorem}

\begin{proof}
  将下确界记为$c>0$。定义保证任意同一个$M\in\mathcal M$均满足
  $J_M(\pi_1)-J_M(\pi_2)\geq c>0$，于是每个模型的排序相同。比较两策略时未切换模型，下确界也无须由某个模型取到，不等式仍成立。
\end{proof}

\begin{example}[不知道共同偏移，仍知道谁更好]
  \label{ex:rl-information-partial-order}
  构造由$u\in[0,10]$参数化的一致模型，并令
  $J_u(\pi_1)=u+1$、$J_u(\pi_2)=u$。两个价值包络$[1,11]$与$[0,10]$大幅重叠，同一模型中的差却恒为1，因此$\pi_1$总较好。若将单独包络端点相减，$1-10=-9$对应为$\pi_1$取$u=0$、为$\pi_2$取$u=10$；这是松的包络差，不是任一共同模型的真实性能差。
\end{example}

这里两候选共享未识别收益$u$，而“多得1”由额外结构假设给定，不能从拟合日志自行宣称。若允许策略只有$\pi_1,\pi_2$，为使选择在每个相容模型中都有依据，可以先各自对$u\in[0,10]$取最坏回报，再在两策略中最大化；最坏价值1和0会选中$\pi_1$。直接计算相对$\pi_2$的共同模型差下界则给出改进1。若另加不得偏离已覆盖历史的约束，且它排除了$\pi_1$的新历史，允许类只剩$\pi_2$，这项已由结构支持的改进也被放弃。保守限制可防止无依据推断，也可能因策略类过窄损失回报；系统的鲁棒优化见第~\ref{chap:rl-trustworthy}~章。

增加同类样本可以缩小采样误差，却不会自动排除总体上仍相容的模型。图~\ref{fig:rl-information-identification-width}~假设总体识别区间$[1,3]$，仅为示意将外层统计半径设为$2/\sqrt n$。样本增加时，外层范围缩到$[1,3]$而非单点。这是区分两种宽度的解析草图，不是算法经验曲线，也不是已经推导的置信保证。

\begin{figure}[htbp]
  \centering
  % Analytic illustration only: identified interval [1,3], outer radius 2/sqrt(n).
\begin{tikzpicture}
  \begin{axis}[
    width=\linewidth,height=61mm,
    xmode=log,log basis x=2,xmin=4,xmax=256,
    ymin=0,ymax=4.35,xtick={4,16,64,256},xticklabels={4,16,64,256},
    ytick={0,1,2,3,4},axis lines=left,
    xlabel={样本量$n$},ylabel={价值（回报）},
    tick label style={font=\figurefont\fontsize{8}{10}\selectfont},
    label style={font=\figurefont\fontsize{8}{10}\selectfont},
    axis line style={BookMuted,line width=.7pt},
    tick style={BookMuted},clip=false
  ]
    \path[fill=BookTeal!8] (axis cs:4,1) rectangle (axis cs:256,3);
    \addplot[BookTeal,dashed,line width=.9pt,domain=4:256] {1};
    \addplot[BookTeal,dashed,line width=.9pt,domain=4:256] {3};
    \addplot[BookBlue,line width=1.1pt,domain=4:256,samples=80]
      {1-2/sqrt(x)};
    \addplot[BookBlue,line width=1.1pt,domain=4:256,samples=80]
      {3+2/sqrt(x)};
    \node[font=\figurefont\fontsize{8}{10}\selectfont,text=BookInk]
      at (axis cs:30,2) {不可消除的识别宽度};
    \draw[{Latex[length=1.6mm]}-{Latex[length=1.6mm]},
      BookAmber,line width=.8pt]
      (axis cs:8,3.04) -- (axis cs:8,3.66);
    \node[anchor=west,font=\figurefont\fontsize{7.5}{9}\selectfont,
      text=BookAmber] at (axis cs:9,4.02) {统计误差示意};
  \end{axis}
\end{tikzpicture}

  \caption{教学识别区间$[1,3]$与示意外层范围$[1-2/\sqrt n,3+2/\sqrt n]$。实线范围随样本量缩小，虚线之间的非零宽度来自允许模型仍有差异；半径公式仅用于解释两种不确定性。}
  \label{fig:rl-information-identification-width}
\end{figure}

% Outline: 4.6.5
\subsection{不可识别时的学习边界与信息补充}
\label{sec:rl-information-unidentifiable}

若一致模型仍支持相反排序，当前目标决策就没有被识别。例~\ref{ex:rl-information-confounding}~的日志在两个模型下似然完全相同，最优动作却分别为$a_0$、$a_1$；缺口是专家选择所依据的未记录类型及其他动作的反事实后果。更细的当前观测编码不会补出这些数据，复杂模型收敛到任一解释也不能证明排序正确。

可以通过收缩目标或补充依据继续推进。缩小$\Pi$只比较日志支持的策略，或将全局最优改为基线附近的有限决策，改变的是要回答的问题。增加已知动力学、行为随机化或代理观测等可核验假设，改变的是允许模型集合。新假设须有外部依据才能排除相容解释，方便拟合的设定不能冒充原日志中的事实。

\begin{table*}[htbp]
  \centering
  \small
  \begin{tabularx}{\linewidth}{p{26mm}XX}
    \toprule
    明确的信息缺口 & 可能补充的信息 & 仍须核验的条件 \\
    \midrule
    较早提示未记录 & 从真实回合开始的长历史或足够记忆摘要 & 提示是否确实保留；部署可用性与回合边界一致 \\
    行为机制不明 & 当时动作概率、策略版本及相关上下文 & 概率条件信息正确；潜在类型上的支持仍可能缺失 \\
    当前状态混叠 & 新传感器、主动辨识动作或代理变量 & 新观测能否区分与控制有关的状态；代理识别条件 \\
    目标行为无覆盖 & 受控随机交互、指定区域的实验 & 查询确实覆盖目标历史动作；代价和安全权限允许 \\
    初始或任务条件不明 & 初始状态、目标及采样分层元数据 & 评价分布与部署分布一致，或有正确校正 \\
    \bottomrule
  \end{tabularx}
  \caption{补充信息要对应具体缺口。记录隐藏类型可以消除部分混杂，却未必补齐每种类型下未尝试的动作；新增字段与新增行为各有不同作用。}
  \label{tab:rl-information-information-repair}
\end{table*}

表~\ref{tab:rl-information-information-repair}~逐项对应了可能补充的信息。单独补记概率未必够：例~\ref{ex:rl-information-confounding}~中专家在每种类型上都确定性行动，即使记录$U$，该类型的另一动作仍无支持。识别固定策略价值，还需在允许类型内受控随机行动，或有依据的跨类型、跨动作结构。更长的同类日志不能替代这些行为或结构信息。

若允许有限交互，可以分别安排状态辨识与策略验证预算。辨识选择能区分相容模型的传感或动作，并在需要时更新表示；验证则固定新策略及输入协议，取得独立性能证据。辨识中的自适应数据不能未经处理便当作独立最终测试。权限不足，或新反馈仍在相容环境中同分布时，无法判断的结论应保留。网络容量可帮助降低表示或拟合误差，却不能从相同输入辨清本来不可区分的环境。

% Outline: 4.7
\section{本章小结}
\label{sec:rl-information-summary}

章首的小车问题现在可以分别回答。速度不可见时，需要由历史保留线索，或用信念与记忆形成决策依据；缺少左转经验时，需要检查收集机制与覆盖，判断左转后果能否从日志确定。Tiger说明模型已知仍可能值得付费辨认当前状态，未覆盖动作反例则说明状态完全可见也仍可能无法识别价值。仅增加记忆不能补齐动作后果，仅补充同类行为也不能找回未记录的线索。

即使识别排除了相容环境给出不同答案的歧义，估计仍需控制有限样本、模型与概率误差。优化又会改变被查询的行为，可能利用估计外推；由价值、行为生成器或轨迹模型提取的执行策略，最终仍须在共同评价条件下比较。行为约束和悲观目标限制的是偏离或乐观预测，其名称与小训练损失不能代替误差依据。因此，可识别、可估计、可优化与可验证只有附带条件的联系，不会自动相互推出。

两类限制叠加时，执行历史、日志字段和收集者信息要共同核对。在完整历史、共同环境与初始条件、正确随机化机制和目标支持下，隐藏状态可从共同反馈项中积分，概率比只含历史条件动作概率。存在未记录行为依据时，相同日志却可能对应相反排序。部分识别仍可能通过共同模型差作出有限决策，排序不可识别时则需收缩结论或补充能区分模型的信息。随策略结论报告执行信息、日志机制、覆盖范围、表示假设、查询预算及误差依据，才能确定它适用于哪里。
